% !TEX program = lualatex
% ---------------------------------------------------------------------------------------------
%  Readback note 4: The main inequality, the bound for the primes, and the zero-free half-plane
%
%  Build:  python3 tools/rb.py note <slug>
%
%  One unit of a Lean artifact, rewritten as ordinary mathematics. The prose is hand-written. Lean names,
%  line ranges and counts come from data.tex; quoted Lean statements are copied from the artifact by the build.
%
%  Rules of the text: standard terminology; every symbol is explained where it first appears and listed again
%  in the notation at the end; no letter has two meanings; a proof is complete or says what it takes for granted;
%  Lean names appear only in the LEAN lines and in the concordance.
%
%  The signature rule (the build enforces it): a function, relation or predicate that the artifact defines is
%  written only through its notation macro, which the build makes from the spec (`notation`) with one slot for
%  each explicit argument of the Lean definition. No argument is fixed and dropped, also not "for brevity" and
%  not for one section. Every argument and every hypothesis of a stated declaration stands inside its
%  statement, or in a Setting that the statement cites by number; only implicit and instance arguments (the
%  ambient objects) may stand in the conventions.
% ---------------------------------------------------------------------------------------------
\documentclass[11pt,a4paper]{article}
\ifdefined\pdfvariable\pdfvariable suppressoptionalinfo \numexpr 1+2+4+8\relax\fi
\usepackage[a4paper,hmargin=35mm,top=22mm,bottom=27mm,footskip=12mm,nohead]{geometry}

\usepackage{fontspec}
\directlua{
  luaotfload.add_fallback("gfallback", {"STIXTwoMath-Regular.otf:mode=node;", "DejaVuSans.ttf:mode=node;"})
  luaotfload.add_fallback("gmonofallback", {"DejaVuSansMono.ttf:mode=node;", "STIXTwoMath-Regular.otf:mode=node;"})
}
\setmainfont{NewCM10}[Extension=.otf,UprightFont=*-Book,ItalicFont=*-BookItalic,
  BoldFont=*-Bold,BoldItalicFont=*-BoldItalic,SmallCapsFont=*-Book,SmallCapsFeatures={Letters=SmallCaps},
  RawFeature={fallback=gfallback}]
\setsansfont{NewCMSans10}[Extension=.otf,UprightFont=*-Book,ItalicFont=*-BookOblique,
  BoldFont=*-Bold,BoldItalicFont=*-BoldOblique,RawFeature={fallback=gfallback}]
\setmonofont{JetBrainsMono}[Extension=.otf,UprightFont=*-Regular,ItalicFont=*-Italic,
  BoldFont=*-SemiBold,BoldItalicFont=*-SemiBoldItalic,Scale=MatchLowercase,
  RawFeature={-calt;-liga;fallback=gmonofallback}]
\usepackage{amsmath}
\usepackage{amsthm}
\usepackage[math-style=TeX]{unicode-math}
\setmathfont{NewCMMath-Book.otf}
\emergencystretch=2.5em\relax

\usepackage{xcolor}
\usepackage{tikz}
\usetikzlibrary{calc,arrows.meta,positioning,backgrounds}
\usepackage[protrusion=true,expansion=true]{microtype}
\usepackage{enumitem}
\usepackage{needspace}
\usepackage{booktabs}
\usepackage{array}
\usepackage{longtable}
\usepackage{fancyhdr}
\usepackage{lastpage}
\usepackage{multicol}
\usepackage[hidelinks,pdfusetitle]{hyperref}

\definecolor{ink}{HTML}{0F172A}
\definecolor{inksoft}{HTML}{475569}
\definecolor{inkfaint}{HTML}{94A3B8}
\definecolor{rulegrey}{HTML}{CBD5E1}
\definecolor{arith}{HTML}{C2410C}
\definecolor{est}{HTML}{7C3AED}
\definecolor{ana}{HTML}{0891B2}
\definecolor{ok}{HTML}{047857}
\definecolor{leancol}{HTML}{3730A3}

\title{The main inequality, the bound for the primes, and the zero-free half-plane}
\author{grh-compact}
\date{}

% ------------------------------------------------------------------------------------- macros
\newcommand{\leanfont}{\ttfamily\fontsize{8.6}{10.5}\selectfont\color{leancol}}
\newcommand{\lean}[1]{\csname LN@#1\endcsname}
\newcommand{\fillin}[1]{{\color{arith}\itshape[#1]}}  % a place of the template that is not written yet; the build reports every one that is left
\newcommand{\leanquote}[1]{\csname LQ@#1\endcsname}
% generated by rb.py note; do not edit
\newcommand{\ConcDecls}{48}
\newcommand{\Pin}{72fead10}
\newcommand{\PinText}{that of commit {\leanfont 72fead10}}
\newcommand{\ArtifactName}{grh-compact}
\expandafter\newcommand\csname LN@star\endcsname{\href{https://github.com/lsernau/grh-lean/blob/72fead102f663d5c7a407964a13dc0cf9d977207/grh-compact/src/New/Main.lean\#L121-L145}{{\leanfont star\_inequality}}}
\expandafter\newcommand\csname LQ@star\endcsname{\leanquoteline{0}{lemma star\_inequality \{P\} \{D : Data P\} : D.m * D.remainder ≤ D.logB \textasciicircum{} 2 * D.budget}}
\expandafter\newcommand\csname LN@primeBound\endcsname{\href{https://github.com/lsernau/grh-lean/blob/72fead102f663d5c7a407964a13dc0cf9d977207/grh-compact/src/New/Main.lean\#L149-L167}{{\leanfont eventually\_primeError\_le\_rpow}}}
\expandafter\newcommand\csname LQ@primeBound\endcsname{\leanquoteline{0}{lemma eventually\_primeError\_le\_rpow (P : Parameters) \{ε : ℝ\} (h : 0 < ε / 6 ∧ ε < 1 / 100) :}\leanquoteline{4}{∀ᶠ x in atTop, primeError P.d P.N₀ x ≤ x \textasciicircum{} (8 * ε + 2⁻¹ - 1)}}
\expandafter\newcommand\csname LN@grh\endcsname{\href{https://github.com/lsernau/grh-lean/blob/72fead102f663d5c7a407964a13dc0cf9d977207/grh-compact/src/New/Main.lean\#L174-L190}{{\leanfont New.\allowbreak{}grh}}}
\expandafter\newcommand\csname LQ@grh\endcsname{\leanquoteline{0}{lemma New.grh : DirichletLNonvanishingRightOfHalf}}
\expandafter\newcommand\csname LN@zerosOnLine\endcsname{\href{https://github.com/lsernau/grh-lean/blob/72fead102f663d5c7a407964a13dc0cf9d977207/grh-compact/src/New/Main.lean\#L194-L196}{{\leanfont New.\allowbreak{}zeros\_on\_line}}}
\expandafter\newcommand\csname LQ@zerosOnLine\endcsname{\leanquoteline{0}{lemma New.zeros\_on\_line (h : DirichletLNonvanishingRightOfHalf) : DirichletLZerosOnLine}}
\expandafter\newcommand\csname LN@zetaZerosOnLine\endcsname{\href{https://github.com/lsernau/grh-lean/blob/72fead102f663d5c7a407964a13dc0cf9d977207/grh-compact/src/New/Main.lean\#L200-L204}{{\leanfont New.\allowbreak{}zeta\_zeros\_on\_line}}}
\expandafter\newcommand\csname LQ@zetaZerosOnLine\endcsname{\leanquoteline{0}{lemma New.zeta\_zeros\_on\_line (h : DirichletLNonvanishingRightOfHalf) :}\leanquoteline{4}{RiemannZetaZerosOnLine}}
\expandafter\newcommand\csname LN@rh\endcsname{\href{https://github.com/lsernau/grh-lean/blob/72fead102f663d5c7a407964a13dc0cf9d977207/grh-compact/src/Statements.lean\#L3-L4}{{\leanfont RiemannZetaZerosOnLine}}}
\expandafter\newcommand\csname LN@halfPlane\endcsname{\href{https://github.com/lsernau/grh-lean/blob/72fead102f663d5c7a407964a13dc0cf9d977207/grh-compact/src/Statements.lean\#L6-L8}{{\leanfont DirichletLNonvanishingRightOfHalf}}}
\expandafter\newcommand\csname LN@strip\endcsname{\href{https://github.com/lsernau/grh-lean/blob/72fead102f663d5c7a407964a13dc0cf9d977207/grh-compact/src/Statements.lean\#L10-L12}{{\leanfont DirichletLZerosOnLine}}}
\expandafter\newcommand\csname LN@params\endcsname{\href{https://github.com/lsernau/grh-lean/blob/72fead102f663d5c7a407964a13dc0cf9d977207/grh-compact/src/New/Data/Different.lean\#L23-L28}{{\leanfont Parameters}}}
\expandafter\newcommand\csname LN@adm\endcsname{\href{https://github.com/lsernau/grh-lean/blob/72fead102f663d5c7a407964a13dc0cf9d977207/grh-compact/src/New/Data/AdmissibleData.lean\#L22-L45}{{\leanfont Data}}}
\expandafter\newcommand\csname LN@cP\endcsname{\href{https://github.com/lsernau/grh-lean/blob/72fead102f663d5c7a407964a13dc0cf9d977207/grh-compact/src/New/Data/AdmissibleData.lean\#L15-L17}{{\leanfont Parameters.\allowbreak{}c}}}
\expandafter\newcommand\csname LN@None\endcsname{\href{https://github.com/lsernau/grh-lean/blob/72fead102f663d5c7a407964a13dc0cf9d977207/grh-compact/src/New/Data/AdmissibleData.lean\#L69-L69}{{\leanfont Data.\allowbreak{}N₁}}}
\expandafter\newcommand\csname LN@ED\endcsname{\href{https://github.com/lsernau/grh-lean/blob/72fead102f663d5c7a407964a13dc0cf9d977207/grh-compact/src/New/Data/AdmissibleData.lean\#L62-L62}{{\leanfont Data.\allowbreak{}E}}}
\expandafter\newcommand\csname LN@aD\endcsname{\href{https://github.com/lsernau/grh-lean/blob/72fead102f663d5c7a407964a13dc0cf9d977207/grh-compact/src/New/Data/AdmissibleData.lean\#L101-L101}{{\leanfont Data.\allowbreak{}a}}}
\expandafter\newcommand\csname LN@fD\endcsname{\href{https://github.com/lsernau/grh-lean/blob/72fead102f663d5c7a407964a13dc0cf9d977207/grh-compact/src/New/Data/AdmissibleData.lean\#L72-L72}{{\leanfont Data.\allowbreak{}f}}}
\expandafter\newcommand\csname LN@tops\endcsname{\href{https://github.com/lsernau/grh-lean/blob/72fead102f663d5c7a407964a13dc0cf9d977207/grh-compact/src/New/Data/AdmissibleData.lean\#L127-L127}{{\leanfont Data.\allowbreak{}tops}}}
\expandafter\newcommand\csname LN@budget\endcsname{\href{https://github.com/lsernau/grh-lean/blob/72fead102f663d5c7a407964a13dc0cf9d977207/grh-compact/src/New/Data/AdmissibleData.lean\#L143-L143}{{\leanfont Data.\allowbreak{}budget}}}
\expandafter\newcommand\csname LN@lam\endcsname{\href{https://github.com/lsernau/grh-lean/blob/72fead102f663d5c7a407964a13dc0cf9d977207/grh-compact/src/New/Endgame/PrimeError.lean\#L25-L25}{{\leanfont lambdaClass}}}
\expandafter\newcommand\csname LN@entireL\endcsname{\href{https://github.com/lsernau/grh-lean/blob/72fead102f663d5c7a407964a13dc0cf9d977207/grh-compact/src/Lib/General/NumberTheory/QuadField.lean\#L31-L31}{{\leanfont entireL}}}
\expandafter\newcommand\csname LN@R\endcsname{\href{https://github.com/lsernau/grh-lean/blob/72fead102f663d5c7a407964a13dc0cf9d977207/grh-compact/src/New/Endgame/PrimeError.lean\#L38-L39}{{\leanfont primeError}}}
\expandafter\newcommand\csname LN@nomW\endcsname{\href{https://github.com/lsernau/grh-lean/blob/72fead102f663d5c7a407964a13dc0cf9d977207/grh-compact/src/New/Data/AdmissibleData.lean\#L134-L134}{{\leanfont Data.\allowbreak{}nominalWeight}}}
\expandafter\newcommand\csname LN@Ask\endcsname{\href{https://github.com/lsernau/grh-lean/blob/72fead102f663d5c7a407964a13dc0cf9d977207/grh-compact/src/New/Data/AdmissibleData.lean\#L138-L139}{{\leanfont Data.\allowbreak{}primeSum}}}
\expandafter\newcommand\csname LN@chowla\endcsname{\href{https://github.com/lsernau/grh-lean/blob/72fead102f663d5c7a407964a13dc0cf9d977207/grh-compact/src/Lib/General/NumberTheory/QuadField.lean\#L35-L35}{{\leanfont chowlaSelbergConst}}}
\expandafter\newcommand\csname LN@mainAsk\endcsname{\href{https://github.com/lsernau/grh-lean/blob/72fead102f663d5c7a407964a13dc0cf9d977207/grh-compact/src/New/Endgame/ShortAndEnough.lean\#L36-L37}{{\leanfont Data.\allowbreak{}mainAsk}}}
\expandafter\newcommand\csname LN@rem\endcsname{\href{https://github.com/lsernau/grh-lean/blob/72fead102f663d5c7a407964a13dc0cf9d977207/grh-compact/src/New/Endgame/ShortAndEnough.lean\#L42-L42}{{\leanfont Data.\allowbreak{}remainder}}}
\expandafter\newcommand\csname LN@rows\endcsname{\href{https://github.com/lsernau/grh-lean/blob/72fead102f663d5c7a407964a13dc0cf9d977207/grh-compact/src/New/Trace/GoodIdempotent.lean\#L118-L118}{{\leanfont Data.\allowbreak{}rows}}}
\expandafter\newcommand\csname LN@rowDim\endcsname{\href{https://github.com/lsernau/grh-lean/blob/72fead102f663d5c7a407964a13dc0cf9d977207/grh-compact/src/New/Trace/GoodIdempotent.lean\#L121-L121}{{\leanfont Data.\allowbreak{}rowDim}}}
\expandafter\newcommand\csname LN@Lift\endcsname{\href{https://github.com/lsernau/grh-lean/blob/72fead102f663d5c7a407964a13dc0cf9d977207/grh-compact/src/New/Rows/RowVolume.lean\#L27-L27}{{\leanfont Data.\allowbreak{}Θ}}}
\expandafter\newcommand\csname LN@covol\endcsname{\href{https://github.com/lsernau/grh-lean/blob/72fead102f663d5c7a407964a13dc0cf9d977207/grh-compact/src/Lib/Adelic/CuspForms/LevelForms.lean\#L206-L206}{{\leanfont LevelPlaces.\allowbreak{}LevelForm.\allowbreak{}covolumeSq}}}
\expandafter\newcommand\csname LN@indexLe\endcsname{\href{https://github.com/lsernau/grh-lean/blob/72fead102f663d5c7a407964a13dc0cf9d977207/grh-compact/src/New/Endgame/ShortAndEnough.lean\#L130-L132}{{\leanfont Data.\allowbreak{}IndexLe}}}
\expandafter\newcommand\csname LN@hmain\endcsname{\href{https://github.com/lsernau/grh-lean/blob/72fead102f663d5c7a407964a13dc0cf9d977207/grh-compact/src/New/Endgame/ShortAndEnough.lean\#L32-L32}{{\leanfont Data.\allowbreak{}mainHeight}}}
\expandafter\newcommand\csname LN@familyAB\endcsname{\href{https://github.com/lsernau/grh-lean/blob/72fead102f663d5c7a407964a13dc0cf9d977207/grh-compact/src/New/Endgame/LongForms.lean\#L78-L92}{{\leanfont Data.\allowbreak{}exists\_family\_covolumeSq\_le}}}
\expandafter\newcommand\csname LQ@familyAB\endcsname{\leanquoteline{0}{lemma Data.exists\_family\_covolumeSq\_le (D : Data P) : ∃ n, ∃ y : Fin n → \_, LinearIndependent ℂ y ∧}\leanquoteline{4}{(1 - 4 / D.b \textasciicircum{} 2) * D.rowDim ≤ n ∧}\leanquoteline{2}{(∃ δ ≤ 5 * D.rowDim * log (1 + D.b \textasciicircum{} 2) * D.budget, D.IndexLe y δ) ∧ ∀ σ,}\leanquoteline{2}{covolumeSq (D.Θ 1 ∘ y) σ ≤ exp (n * (D.mainHeight - D.m * D.remainder + D.logB * D.budget))}}
\expandafter\newcommand\csname LN@longC\endcsname{\href{https://github.com/lsernau/grh-lean/blob/72fead102f663d5c7a407964a13dc0cf9d977207/grh-compact/src/New/Endgame/LongForms.lean\#L63-L71}{{\leanfont Data.\allowbreak{}integral\_forms\_long}}}
\expandafter\newcommand\csname LQ@longC\endcsname{\leanquoteline{0}{lemma Data.integral\_forms\_long : D.rowDim * (D.mainHeight - 3 * (D.logB * D.budget)) ≤}\leanquoteline{2}{V + 2 * δ}}
\expandafter\newcommand\csname LN@crude\endcsname{\href{https://github.com/lsernau/grh-lean/blob/72fead102f663d5c7a407964a13dc0cf9d977207/grh-compact/src/New/Endgame/ShortAndEnough.lean\#L74-L82}{{\leanfont Data.\allowbreak{}abs\_mainHeight\_add\_remainder\_le}}}
\expandafter\newcommand\csname LN@rowsNe\endcsname{\href{https://github.com/lsernau/grh-lean/blob/72fead102f663d5c7a407964a13dc0cf9d977207/grh-compact/src/New/Filling/FillingTheorem.lean\#L45-L47}{{\leanfont Data.\allowbreak{}finrank\_rows\_ne\_zero}}}
\expandafter\newcommand\csname LN@cGe\endcsname{\href{https://github.com/lsernau/grh-lean/blob/72fead102f663d5c7a407964a13dc0cf9d977207/grh-compact/src/New/Data/Magnitudes.lean\#L21-L22}{{\leanfont Data.\allowbreak{}c\_ge}}}
\expandafter\newcommand\csname LN@cardPlaces\endcsname{\href{https://github.com/lsernau/grh-lean/blob/72fead102f663d5c7a407964a13dc0cf9d977207/grh-compact/src/New/Data/Different.lean\#L168-L170}{{\leanfont Parameters.\allowbreak{}card\_infinitePlace\_E}}}
\expandafter\newcommand\csname LN@cardTops\endcsname{\href{https://github.com/lsernau/grh-lean/blob/72fead102f663d5c7a407964a13dc0cf9d977207/grh-compact/src/New/Data/Different.lean\#L174-L176}{{\leanfont Parameters.\allowbreak{}card\_placeClass\_eq}}}
\expandafter\newcommand\csname LN@remEq\endcsname{\href{https://github.com/lsernau/grh-lean/blob/72fead102f663d5c7a407964a13dc0cf9d977207/grh-compact/src/New/Endgame/ShortAndEnough.lean\#L86-L112}{{\leanfont Data.\allowbreak{}remainder\_eq}}}
\expandafter\newcommand\csname LN@existsData\endcsname{\href{https://github.com/lsernau/grh-lean/blob/72fead102f663d5c7a407964a13dc0cf9d977207/grh-compact/src/New/Data/DataExist.lean\#L49-L108}{{\leanfont exists\_data}}}
\expandafter\newcommand\csname LN@isPolar\endcsname{\href{https://github.com/lsernau/grh-lean/blob/72fead102f663d5c7a407964a13dc0cf9d977207/grh-compact/src/Lib/Analytic/Landau/PolarLandau.lean\#L34-L36}{{\leanfont IsPolar}}}
\expandafter\newcommand\csname LN@polarR\endcsname{\href{https://github.com/lsernau/grh-lean/blob/72fead102f663d5c7a407964a13dc0cf9d977207/grh-compact/src/New/Endgame/PrimeError.lean\#L69-L111}{{\leanfont isPolar\_primeError}}}
\expandafter\newcommand\csname LN@window\endcsname{\href{https://github.com/lsernau/grh-lean/blob/72fead102f663d5c7a407964a13dc0cf9d977207/grh-compact/src/New/Endgame/OneSidedCriterion.lean\#L73-L99}{{\leanfont IsPolar.\allowbreak{}exists\_window\_sub\_mul\_le}}}
\expandafter\newcommand\csname LN@criterion\endcsname{\href{https://github.com/lsernau/grh-lean/blob/72fead102f663d5c7a407964a13dc0cf9d977207/grh-compact/src/New/Endgame/OneSidedCriterion.lean\#L28-L36}{{\leanfont entireL\_ne\_zero\_of\_primeError\_le}}}
\expandafter\newcommand\csname LN@symm\endcsname{\href{https://github.com/lsernau/grh-lean/blob/72fead102f663d5c7a407964a13dc0cf9d977207/grh-compact/src/Lib/General/NumberTheory/ZeroSymmetry.lean\#L27-L33}{{\leanfont LFunction\_one\_sub\_eq\_zero}}}
\expandafter\newcommand\csname LN@heeg\endcsname{\href{https://github.com/lsernau/grh-lean/blob/72fead102f663d5c7a407964a13dc0cf9d977207/grh-compact/src/New/Data/Different.lean\#L36-L37}{{\leanfont Parameters.\allowbreak{}instHeegD}}}
\expandafter\newcommand\csname LN@coprime\endcsname{\href{https://github.com/lsernau/grh-lean/blob/72fead102f663d5c7a407964a13dc0cf9d977207/grh-compact/src/New/Data/Different.lean\#L74-L75}{{\leanfont Parameters.\allowbreak{}coprime\_N₁\_d}}}
\expandafter\newcommand\csname LN@sizeLe\endcsname{\href{https://github.com/lsernau/grh-lean/blob/72fead102f663d5c7a407964a13dc0cf9d977207/grh-compact/src/New/Data/Magnitudes.lean\#L121-L121}{{\leanfont Data.\allowbreak{}le\_of\_hasLogBound}}}
\newcommand{\ConcordanceRows}{\crow{star}{\href{https://github.com/lsernau/grh-lean/blob/72fead102f663d5c7a407964a13dc0cf9d977207/grh-compact/src/New/Main.lean\#L121-L145}{{\leanfont star\_inequality}}}{Main.lean}{121--145}
\crow{primeBound}{\href{https://github.com/lsernau/grh-lean/blob/72fead102f663d5c7a407964a13dc0cf9d977207/grh-compact/src/New/Main.lean\#L149-L167}{{\leanfont eventually\_primeError\_le\_rpow}}}{Main.lean}{149--167}
\crow{grh}{\href{https://github.com/lsernau/grh-lean/blob/72fead102f663d5c7a407964a13dc0cf9d977207/grh-compact/src/New/Main.lean\#L174-L190}{{\leanfont New.\allowbreak{}grh}}}{Main.lean}{174--190}
\crow{zerosOnLine}{\href{https://github.com/lsernau/grh-lean/blob/72fead102f663d5c7a407964a13dc0cf9d977207/grh-compact/src/New/Main.lean\#L194-L196}{{\leanfont New.\allowbreak{}zeros\_on\_line}}}{Main.lean}{194--196}
\crow{zetaZerosOnLine}{\href{https://github.com/lsernau/grh-lean/blob/72fead102f663d5c7a407964a13dc0cf9d977207/grh-compact/src/New/Main.lean\#L200-L204}{{\leanfont New.\allowbreak{}zeta\_zeros\_on\_line}}}{Main.lean}{200--204}
\crow{rh}{\href{https://github.com/lsernau/grh-lean/blob/72fead102f663d5c7a407964a13dc0cf9d977207/grh-compact/src/Statements.lean\#L3-L4}{{\leanfont RiemannZetaZerosOnLine}}}{Statements.lean}{3--4}
\crow{halfPlane}{\href{https://github.com/lsernau/grh-lean/blob/72fead102f663d5c7a407964a13dc0cf9d977207/grh-compact/src/Statements.lean\#L6-L8}{{\leanfont DirichletLNonvanishingRightOfHalf}}}{Statements.lean}{6--8}
\crow{strip}{\href{https://github.com/lsernau/grh-lean/blob/72fead102f663d5c7a407964a13dc0cf9d977207/grh-compact/src/Statements.lean\#L10-L12}{{\leanfont DirichletLZerosOnLine}}}{Statements.lean}{10--12}
\crow{params}{\href{https://github.com/lsernau/grh-lean/blob/72fead102f663d5c7a407964a13dc0cf9d977207/grh-compact/src/New/Data/Different.lean\#L23-L28}{{\leanfont Parameters}}}{Different.lean}{23--28}
\crow{adm}{\href{https://github.com/lsernau/grh-lean/blob/72fead102f663d5c7a407964a13dc0cf9d977207/grh-compact/src/New/Data/AdmissibleData.lean\#L22-L45}{{\leanfont Data}}}{AdmissibleData.lean}{22--45}
\crow{cP}{\href{https://github.com/lsernau/grh-lean/blob/72fead102f663d5c7a407964a13dc0cf9d977207/grh-compact/src/New/Data/AdmissibleData.lean\#L15-L17}{{\leanfont Parameters.\allowbreak{}c}}}{AdmissibleData.lean}{15--17}
\crow{None}{\href{https://github.com/lsernau/grh-lean/blob/72fead102f663d5c7a407964a13dc0cf9d977207/grh-compact/src/New/Data/AdmissibleData.lean\#L69-L69}{{\leanfont Data.\allowbreak{}N₁}}}{AdmissibleData.lean}{69--69}
\crow{ED}{\href{https://github.com/lsernau/grh-lean/blob/72fead102f663d5c7a407964a13dc0cf9d977207/grh-compact/src/New/Data/AdmissibleData.lean\#L62-L62}{{\leanfont Data.\allowbreak{}E}}}{AdmissibleData.lean}{62--62}
\crow{aD}{\href{https://github.com/lsernau/grh-lean/blob/72fead102f663d5c7a407964a13dc0cf9d977207/grh-compact/src/New/Data/AdmissibleData.lean\#L101-L101}{{\leanfont Data.\allowbreak{}a}}}{AdmissibleData.lean}{101--101}
\crow{fD}{\href{https://github.com/lsernau/grh-lean/blob/72fead102f663d5c7a407964a13dc0cf9d977207/grh-compact/src/New/Data/AdmissibleData.lean\#L72-L72}{{\leanfont Data.\allowbreak{}f}}}{AdmissibleData.lean}{72--72}
\crow{tops}{\href{https://github.com/lsernau/grh-lean/blob/72fead102f663d5c7a407964a13dc0cf9d977207/grh-compact/src/New/Data/AdmissibleData.lean\#L127-L127}{{\leanfont Data.\allowbreak{}tops}}}{AdmissibleData.lean}{127--127}
\crow{budget}{\href{https://github.com/lsernau/grh-lean/blob/72fead102f663d5c7a407964a13dc0cf9d977207/grh-compact/src/New/Data/AdmissibleData.lean\#L143-L143}{{\leanfont Data.\allowbreak{}budget}}}{AdmissibleData.lean}{143--143}
\crow{lam}{\href{https://github.com/lsernau/grh-lean/blob/72fead102f663d5c7a407964a13dc0cf9d977207/grh-compact/src/New/Endgame/PrimeError.lean\#L25-L25}{{\leanfont lambdaClass}}}{PrimeError.lean}{25--25}
\crow{entireL}{\href{https://github.com/lsernau/grh-lean/blob/72fead102f663d5c7a407964a13dc0cf9d977207/grh-compact/src/Lib/General/NumberTheory/QuadField.lean\#L31-L31}{{\leanfont entireL}}}{QuadField.lean}{31--31}
\crow{R}{\href{https://github.com/lsernau/grh-lean/blob/72fead102f663d5c7a407964a13dc0cf9d977207/grh-compact/src/New/Endgame/PrimeError.lean\#L38-L39}{{\leanfont primeError}}}{PrimeError.lean}{38--39}
\crow{nomW}{\href{https://github.com/lsernau/grh-lean/blob/72fead102f663d5c7a407964a13dc0cf9d977207/grh-compact/src/New/Data/AdmissibleData.lean\#L134-L134}{{\leanfont Data.\allowbreak{}nominalWeight}}}{AdmissibleData.lean}{134--134}
\crow{Ask}{\href{https://github.com/lsernau/grh-lean/blob/72fead102f663d5c7a407964a13dc0cf9d977207/grh-compact/src/New/Data/AdmissibleData.lean\#L138-L139}{{\leanfont Data.\allowbreak{}primeSum}}}{AdmissibleData.lean}{138--139}
\crow{chowla}{\href{https://github.com/lsernau/grh-lean/blob/72fead102f663d5c7a407964a13dc0cf9d977207/grh-compact/src/Lib/General/NumberTheory/QuadField.lean\#L35-L35}{{\leanfont chowlaSelbergConst}}}{QuadField.lean}{35--35}
\crow{mainAsk}{\href{https://github.com/lsernau/grh-lean/blob/72fead102f663d5c7a407964a13dc0cf9d977207/grh-compact/src/New/Endgame/ShortAndEnough.lean\#L36-L37}{{\leanfont Data.\allowbreak{}mainAsk}}}{ShortAndEnough.lean}{36--37}
\crow{rem}{\href{https://github.com/lsernau/grh-lean/blob/72fead102f663d5c7a407964a13dc0cf9d977207/grh-compact/src/New/Endgame/ShortAndEnough.lean\#L42-L42}{{\leanfont Data.\allowbreak{}remainder}}}{ShortAndEnough.lean}{42--42}
\crow{rows}{\href{https://github.com/lsernau/grh-lean/blob/72fead102f663d5c7a407964a13dc0cf9d977207/grh-compact/src/New/Trace/GoodIdempotent.lean\#L118-L118}{{\leanfont Data.\allowbreak{}rows}}}{GoodIdempotent.lean}{118--118}
\crow{rowDim}{\href{https://github.com/lsernau/grh-lean/blob/72fead102f663d5c7a407964a13dc0cf9d977207/grh-compact/src/New/Trace/GoodIdempotent.lean\#L121-L121}{{\leanfont Data.\allowbreak{}rowDim}}}{GoodIdempotent.lean}{121--121}
\crow{Lift}{\href{https://github.com/lsernau/grh-lean/blob/72fead102f663d5c7a407964a13dc0cf9d977207/grh-compact/src/New/Rows/RowVolume.lean\#L27-L27}{{\leanfont Data.\allowbreak{}Θ}}}{RowVolume.lean}{27--27}
\crow{covol}{\href{https://github.com/lsernau/grh-lean/blob/72fead102f663d5c7a407964a13dc0cf9d977207/grh-compact/src/Lib/Adelic/CuspForms/LevelForms.lean\#L206-L206}{{\leanfont LevelPlaces.\allowbreak{}LevelForm.\allowbreak{}covolumeSq}}}{LevelForms.lean}{206--206}
\crow{indexLe}{\href{https://github.com/lsernau/grh-lean/blob/72fead102f663d5c7a407964a13dc0cf9d977207/grh-compact/src/New/Endgame/ShortAndEnough.lean\#L130-L132}{{\leanfont Data.\allowbreak{}IndexLe}}}{ShortAndEnough.lean}{130--132}
\crow{hmain}{\href{https://github.com/lsernau/grh-lean/blob/72fead102f663d5c7a407964a13dc0cf9d977207/grh-compact/src/New/Endgame/ShortAndEnough.lean\#L32-L32}{{\leanfont Data.\allowbreak{}mainHeight}}}{ShortAndEnough.lean}{32--32}
\crow{familyAB}{\href{https://github.com/lsernau/grh-lean/blob/72fead102f663d5c7a407964a13dc0cf9d977207/grh-compact/src/New/Endgame/LongForms.lean\#L78-L92}{{\leanfont Data.\allowbreak{}exists\_family\_covolumeSq\_le}}}{LongForms.lean}{78--92}
\crow{longC}{\href{https://github.com/lsernau/grh-lean/blob/72fead102f663d5c7a407964a13dc0cf9d977207/grh-compact/src/New/Endgame/LongForms.lean\#L63-L71}{{\leanfont Data.\allowbreak{}integral\_forms\_long}}}{LongForms.lean}{63--71}
\crow{crude}{\href{https://github.com/lsernau/grh-lean/blob/72fead102f663d5c7a407964a13dc0cf9d977207/grh-compact/src/New/Endgame/ShortAndEnough.lean\#L74-L82}{{\leanfont Data.\allowbreak{}abs\_mainHeight\_add\_remainder\_le}}}{ShortAndEnough.lean}{74--82}
\crow{rowsNe}{\href{https://github.com/lsernau/grh-lean/blob/72fead102f663d5c7a407964a13dc0cf9d977207/grh-compact/src/New/Filling/FillingTheorem.lean\#L45-L47}{{\leanfont Data.\allowbreak{}finrank\_rows\_ne\_zero}}}{FillingTheorem.lean}{45--47}
\crow{cGe}{\href{https://github.com/lsernau/grh-lean/blob/72fead102f663d5c7a407964a13dc0cf9d977207/grh-compact/src/New/Data/Magnitudes.lean\#L21-L22}{{\leanfont Data.\allowbreak{}c\_ge}}}{Magnitudes.lean}{21--22}
\crow{cardPlaces}{\href{https://github.com/lsernau/grh-lean/blob/72fead102f663d5c7a407964a13dc0cf9d977207/grh-compact/src/New/Data/Different.lean\#L168-L170}{{\leanfont Parameters.\allowbreak{}card\_infinitePlace\_E}}}{Different.lean}{168--170}
\crow{cardTops}{\href{https://github.com/lsernau/grh-lean/blob/72fead102f663d5c7a407964a13dc0cf9d977207/grh-compact/src/New/Data/Different.lean\#L174-L176}{{\leanfont Parameters.\allowbreak{}card\_placeClass\_eq}}}{Different.lean}{174--176}
\crow{remEq}{\href{https://github.com/lsernau/grh-lean/blob/72fead102f663d5c7a407964a13dc0cf9d977207/grh-compact/src/New/Endgame/ShortAndEnough.lean\#L86-L112}{{\leanfont Data.\allowbreak{}remainder\_eq}}}{ShortAndEnough.lean}{86--112}
\crow{existsData}{\href{https://github.com/lsernau/grh-lean/blob/72fead102f663d5c7a407964a13dc0cf9d977207/grh-compact/src/New/Data/DataExist.lean\#L49-L108}{{\leanfont exists\_data}}}{DataExist.lean}{49--108}
\crow{isPolar}{\href{https://github.com/lsernau/grh-lean/blob/72fead102f663d5c7a407964a13dc0cf9d977207/grh-compact/src/Lib/Analytic/Landau/PolarLandau.lean\#L34-L36}{{\leanfont IsPolar}}}{PolarLandau.lean}{34--36}
\crow{polarR}{\href{https://github.com/lsernau/grh-lean/blob/72fead102f663d5c7a407964a13dc0cf9d977207/grh-compact/src/New/Endgame/PrimeError.lean\#L69-L111}{{\leanfont isPolar\_primeError}}}{PrimeError.lean}{69--111}
\crow{window}{\href{https://github.com/lsernau/grh-lean/blob/72fead102f663d5c7a407964a13dc0cf9d977207/grh-compact/src/New/Endgame/OneSidedCriterion.lean\#L73-L99}{{\leanfont IsPolar.\allowbreak{}exists\_window\_sub\_mul\_le}}}{OneSidedCriterion.lean}{73--99}
\crow{criterion}{\href{https://github.com/lsernau/grh-lean/blob/72fead102f663d5c7a407964a13dc0cf9d977207/grh-compact/src/New/Endgame/OneSidedCriterion.lean\#L28-L36}{{\leanfont entireL\_ne\_zero\_of\_primeError\_le}}}{OneSidedCriterion.lean}{28--36}
\crow{symm}{\href{https://github.com/lsernau/grh-lean/blob/72fead102f663d5c7a407964a13dc0cf9d977207/grh-compact/src/Lib/General/NumberTheory/ZeroSymmetry.lean\#L27-L33}{{\leanfont LFunction\_one\_sub\_eq\_zero}}}{ZeroSymmetry.lean}{27--33}
\crow{heeg}{\href{https://github.com/lsernau/grh-lean/blob/72fead102f663d5c7a407964a13dc0cf9d977207/grh-compact/src/New/Data/Different.lean\#L36-L37}{{\leanfont Parameters.\allowbreak{}instHeegD}}}{Different.lean}{36--37}
\crow{coprime}{\href{https://github.com/lsernau/grh-lean/blob/72fead102f663d5c7a407964a13dc0cf9d977207/grh-compact/src/New/Data/Different.lean\#L74-L75}{{\leanfont Parameters.\allowbreak{}coprime\_N₁\_d}}}{Different.lean}{74--75}
\crow{sizeLe}{\href{https://github.com/lsernau/grh-lean/blob/72fead102f663d5c7a407964a13dc0cf9d977207/grh-compact/src/New/Data/Magnitudes.lean\#L121-L121}{{\leanfont Data.\allowbreak{}le\_of\_hasLogBound}}}{Magnitudes.lean}{121--121}}
% the notations of the defined functions: one slot for each explicit argument of the Lean declaration
\newcommand{\cP}[1]{\symit{C}(#1)}
\newcommand{\None}[1]{N_1(#1)}
\newcommand{\ED}[1]{E(#1)}
\newcommand{\aD}[1]{\symit{a}(#1)}
\newcommand{\fD}[1]{\symit{f}(#1)}
\newcommand{\tops}[1]{\mathrm{tops}(#1)}
\newcommand{\bud}[1]{\mathrm{budget}(#1)}
\newcommand{\lam}[3]{\lambda_{#1,#2}(#3)}
\newcommand{\entireL}[1]{\tilde L(\,\cdot\,,#1)}
\newcommand{\entireLs}[1]{\tilde L(s,#1)}
\newcommand{\entireLrho}[1]{\tilde L(\rho,#1)}
\newcommand{\entireLone}[1]{\tilde L(1,#1)}
\newcommand{\R}[3]{R_{#1,#2}(#3)}
\newcommand{\nomW}[2]{w_{#1}(#2)}
\newcommand{\Ask}[3]{\mathrm{Ask}_{#1}(#2,#3)}
\newcommand{\chowla}[1]{\operatorname{cs}(#1)}
\newcommand{\mainAsk}[1]{\mathrm{Ask}^{\mathrm{main}}(#1)}
\newcommand{\rem}[1]{\mathrm{rem}(#1)}
\newcommand{\rows}[1]{\mathrm{Row}(#1)}
\newcommand{\rowDim}[1]{\dim\mathrm{Row}(#1)}
\newcommand{\Lift}[2]{\Theta_{#1,#2}}
\newcommand{\covol}[2]{\operatorname{covol}_{#2}^{2}(#1)}
\newcommand{\indexLe}[3]{#2\text{ have }#1\text{-index at most }#3}
\newcommand{\hmain}[1]{h_{\mathrm{main}}(#1)}
\newcommand{\polar}[2]{\operatorname{Polar}(#1,#2)}
\newcommand{\polarWords}[2]{$#1$ is polar with residue function~$#2$}
\newcommand{\UnitMembers}{5}
\newcommand{\UnitInNote}{5}
\newcommand{\UnitProofs}{5}
\newcommand{\UnitSteps}{27}
\newcommand{\UnitAllSteps}{34}


\linespread{1.08}                       % 11 pt on 14.7 pt
\setlength{\parindent}{0pt}
\setlength{\parskip}{5.5pt plus 1.5pt minus 0.5pt}
\setlength{\abovedisplayskip}{7.5pt plus 2pt minus 3pt}
\setlength{\belowdisplayskip}{7.5pt plus 2pt minus 3pt}
\setlength{\abovedisplayshortskip}{3pt plus 2pt}
\setlength{\belowdisplayshortskip}{6pt plus 2pt minus 2pt}
\emergencystretch=1.5em
\clubpenalty=10000 \widowpenalty=10000 \displaywidowpenalty=10000
\relpenalty=7000 \binoppenalty=9000
\raggedbottom
\pagestyle{fancy}\fancyhf{}\renewcommand{\headrulewidth}{0pt}
\fancyfoot[C]{\sffamily\fontsize{7.6}{9}\selectfont\color{inkfaint}\thepage\,/\,\pageref*{LastPage}}
\setlist{topsep=2pt,partopsep=0pt,parsep=0pt,itemsep=4pt,leftmargin=2.1em,labelsep=0.55em}

\newcommand{\Q}{\mathbb{Q}}
\newcommand{\Z}{\mathbb{Z}}
\newcommand{\F}{\mathbb{F}}
\newcommand{\len}{\operatorname{len}}
\newcommand{\cok}{\operatorname{cok}}
\newcommand{\End}{\operatorname{End}}
\newcommand{\Hom}{\operatorname{Hom}}
\newcommand{\Gal}{\operatorname{Gal}}
\newcommand{\diag}{\operatorname{diag}}
\newcommand{\kw}[1]{\emph{#1}}
\newcommand{\m}[1]{\mbox{$#1$}}        % a short formula that must not be broken across lines
\newcommand{\head}[1]{\par\addvspace{17pt plus 4pt minus 2pt}\Needspace{6\baselineskip}%
  \noindent{\fontsize{13.2}{16}\selectfont\bfseries\color{ink}#1}\par\nopagebreak\vspace{-1pt}}
\newcommand{\sect}[2]{\refstepcounter{section}\label{#2}\par\addvspace{19pt plus 4pt minus 2pt}\Needspace{7\baselineskip}%
  \noindent\llap{{\fontsize{13.2}{16}\selectfont\bfseries\color{inkfaint}\thesection}\hspace{9pt}}%
  {\fontsize{13.2}{16}\selectfont\bfseries\color{ink}#1}\par\nopagebreak\vspace{-1pt}}
\newcommand{\subhead}[1]{\par\addvspace{9pt plus 2pt}\Needspace{4\baselineskip}\noindent\textbf{#1}\enspace\ignorespaces}
% where a statement lives in the Lean archive
\newcommand{\leantag}{{\sffamily\bfseries\fontsize{6.3}{8}\selectfont\color{leancol}\addfontfeatures{LetterSpace=7}LEAN}}
\newcommand{\sep}{\unskip\nobreak\hspace{5pt}{\color{inkfaint}\textperiodcentered}\hspace{5pt plus 2pt}\ignorespaces}
\newcommand{\leanline}[1]{\par\nopagebreak\vspace{-2.5pt}{\fontsize{8.9}{12.4}\selectfont\color{inksoft}\raggedright
  \noindent\llap{\leantag\hspace{9pt}}#1\par}\vspace{1pt}}
% a displayed statement in a tinted box: #1 colour, #2 text
\newcommand{\statement}[2]{\par\addvspace{7pt}\begingroup\setlength{\fboxsep}{8.5pt}\setlength{\fboxrule}{0.6pt}%
  \noindent\fcolorbox{#1!30!white}{#1!3!white}{\parbox{\dimexpr\linewidth-2\fboxsep-2\fboxrule\relax}{#2}}%
  \endgroup\par\addvspace{7pt}}
\newcommand{\nt}[2]{\par\noindent\hangindent=1.1em\hangafter=1 #1\enspace{\color{inksoft}#2}\par\vspace{1pt}}
\newcommand{\legdot}[1]{\tikz[baseline=-0.62ex]\fill[#1] (0,0) circle (2.5pt);\hspace{3pt}}
\newcommand{\crow}[4]{\csname CS@#1\endcsname & #2 & \textcolor{inksoft}{#3} & \textcolor{inksoft}{#4}\\}
\newcommand{\leanquoteline}[2]{\par\noindent\hangindent=\dimexpr#1\fontdimen2\font+2.2em\relax\hangafter=1
  \hspace*{#1\fontdimen2\font}#2}
\newcommand{\quoted}[1]{\par\vspace{3pt}{\ttfamily\fontsize{7.5}{10.4}\selectfont\color{leancol}\raggedright
  \setlength{\parskip}{0pt}\leanquote{#1}\par}\vspace{4pt}}
\newcommand{\bx}[3]{{\sffamily\bfseries\fontsize{7.3}{9.2}\selectfont\color{#1!82!black}#2}\\[1.5pt]#3}
\newcommand{\cs}[2]{\expandafter\newcommand\csname CS@#1\endcsname{#2}}

% theorems
\newtheoremstyle{plainx}{7pt plus 2pt}{7pt plus 2pt}{\itshape}{}{\bfseries}{.}{0.5em}{}
\newtheoremstyle{uprightx}{7pt plus 2pt}{7pt plus 2pt}{}{}{\bfseries}{.}{0.5em}{}
\theoremstyle{plainx}
\newtheorem{theorem}{Theorem}[section]
\newtheorem{lemma}[theorem]{Lemma}
\newtheorem{proposition}[theorem]{Proposition}
\newtheorem{corollary}[theorem]{Corollary}
\theoremstyle{uprightx}
\newtheorem{definition}[theorem]{Definition}
\newtheorem{setting}[theorem]{Setting}          % hypotheses shared by several statements; each of them cites it by number
\newtheorem{remark}[theorem]{Remark}
\newtheorem{example}[theorem]{Example}
\renewcommand{\qedsymbol}{{\setlength{\fboxsep}{0pt}\setlength{\fboxrule}{0.5pt}\fbox{\rule{0pt}{0.46em}\hspace{0.46em}}}}
% what holds throughout the note: the ambient objects only (implicit and instance arguments), never a hypothesis
\newenvironment{conventions}{\par\addvspace{5pt}\noindent\textbf{Conventions.}\enspace\ignorespaces}{\par\addvspace{3pt}}
\makeatletter
\renewenvironment{proof}[1][Proof]{\par\pushQED{\qed}\normalfont\topsep0pt\partopsep0pt
  \trivlist\item[\hskip\labelsep\itshape #1.]\ignorespaces}{\popQED\endtrivlist\@endpefalse}
\makeatother

% where each declaration of the concordance is treated: \cs{key}{place}, the place naming a \label or a lettered item
\cs{star}{Theorem 1}
\cs{primeBound}{Theorem 2}
\cs{grh}{Theorem 3}
\cs{zerosOnLine}{Theorem 4}
\cs{zetaZerosOnLine}{Theorem 5}
\cs{rh}{Definition~\ref{d:statements}}
\cs{halfPlane}{Definition~\ref{d:statements}}
\cs{strip}{Definition~\ref{d:statements}}
\cs{params}{Definition~\ref{d:params}}
\cs{adm}{Definition~\ref{d:data}}
\cs{cP}{Definition~\ref{d:data}}
\cs{None}{Definition~\ref{d:data}}
\cs{ED}{Definition~\ref{d:data}}
\cs{aD}{Definition~\ref{d:derived}}
\cs{fD}{Definition~\ref{d:derived}}
\cs{tops}{Definition~\ref{d:derived}}
\cs{budget}{Definition~\ref{d:derived}}
\cs{lam}{Definition~\ref{d:R}}
\cs{entireL}{Definition~\ref{d:R}}
\cs{R}{Definition~\ref{d:R}}
\cs{nomW}{Definition~\ref{d:rem}}
\cs{Ask}{Definition~\ref{d:rem}}
\cs{chowla}{Definition~\ref{d:rem}}
\cs{mainAsk}{Definition~\ref{d:rem}}
\cs{rem}{Definition~\ref{d:rem}}
\cs{rows}{Definition~\ref{d:later}}
\cs{rowDim}{Definition~\ref{d:later}}
\cs{Lift}{Definition~\ref{d:later}}
\cs{covol}{Definition~\ref{d:later}}
\cs{indexLe}{Definition~\ref{d:later}}
\cs{hmain}{Definition~\ref{d:later}}
\cs{familyAB}{Input 1}
\cs{longC}{Input 2}
\cs{crude}{Input 3}
\cs{rowsNe}{Input 4}
\cs{cGe}{Input 5}
\cs{cardPlaces}{Input 5}
\cs{cardTops}{Input 5}
\cs{remEq}{Input 6}
\cs{existsData}{Input 7}
\cs{isPolar}{Definition~\ref{d:polar}}
\cs{polarR}{Input 8}
\cs{window}{Input 9}
\cs{criterion}{Input 10}
\cs{symm}{Input 11}
\cs{heeg}{Lemma~\ref{l:params}}
\cs{coprime}{Lemma~\ref{l:params}}
\cs{sizeLe}{Lemma~\ref{l:sizes}}
\newcommand{\CC}{\mathbb{C}}
\newcommand{\RR}{\mathbb{R}}
\newcommand{\QQ}{\mathbb{Q}}
\newcommand{\ZZ}{\mathbb{Z}}
\newcommand{\re}{\operatorname{Re}}
\newcommand{\cl}[2]{\operatorname{cl}_{#1}(#2)}
\newcommand{\thick}[1]{\mathrm{thick}(#1)}
\newcommand{\kap}[2]{\kappa_{#1,#2}}
\newcommand{\res}[3]{\operatorname{res}_{#1,#2}(#3)}
\newcommand{\ord}{\operatorname{ord}}
\newcommand{\cE}{\mathcal{E}}
\newcommand{\scx}[1]{x(#1)}
\newcommand{\scxo}[1]{x_1(#1)}
\newcommand{\scb}[1]{b(#1)}
\newcommand{\Dm}[1]{m(#1)}
\newcommand{\Dr}[1]{r(#1)}
\newcommand{\hA}{h_{\mathrm{A}}}
\newcommand{\mlib}[1]{\leavevmode{\leanfont #1}}
\newcommand{\mrow}[2]{#1 & \mlib{#2}\\[1.5pt]}

\begin{document}
\color{ink}

{\raggedright
{\sffamily\bfseries\fontsize{7.2}{9}\selectfont\color{leancol}\addfontfeatures{LetterSpace=8}READBACK NOTE 4\par}
\vspace{5pt}
{\fontsize{21.5}{25.5}\selectfont\bfseries The main inequality, the bound for the primes, and the zero-free
half-plane\par}
\vspace{5pt}
{\fontsize{12.4}{16}\selectfont\itshape\color{inksoft} One unit of the Lean artifact \ArtifactName, read back into
mathematics\par}}
\vspace{8pt}
{\color{rulegrey}\hrule height 0.6pt}
\vspace{7pt}
\thispagestyle{empty}

\textbf{Summary.} This unit is the top of the argument. Its first result, the main inequality, bounds a
remainder of weighted sums over primes: $\Dm{D}\cdot\rem{D}\le(\log \scb{D})^2\cdot\bud{D}$ for all admissible data~$D$ of
the proof. It comes from bounding one covolume from two sides: Theorems A and~B of the argument give many
independent automorphic forms with short theta lifts, Theorem~C says that such lifts cannot be that short, and in the comparison the main amounts cancel. Letting the data grow, with the second scale at a good point of a long
window, one gets the one-sided bound $\R{d}{N_0}{x}\le x^{8\varepsilon-1/2}$ for the remainder of a smoothed sum over the
prime powers $\nu\equiv1$ modulo~$N_0$, weighted by $\Lambda(\nu)\bigl(1-\chi_K(\nu)\bigr)$ (Definition~\ref{d:R}); the saving $x^{-1/2}$ is the reciprocal of the
number~$\Dm{D}$ of lifts. A criterion of Landau turns this bound, for all primes~$d$ of one residue class, into the
absence of zeros of Dirichlet $L$-functions in the half-plane $\re s>\tfrac12$, and the functional equations put
the nontrivial zeros on the critical line.

\statement{leancol}{\fontsize{9.6}{13.2}\selectfont\textbf{What this note is.} The artifact \ArtifactName{} is a
body of Lean code: the compact edition of a formal proof, due to L.~Alpöge with the AI system Claude, whose final
theorems are the Riemann hypothesis and its generalisation to Dirichlet $L$-functions. That headline is a claim.
The authors report that every file compiles with Lean v4.35.0 against a pinned version of the library Mathlib and
that Lean's kernel accepted the proofs; an independent build is recorded in the handout's record of checks when it
reports. No mathematician has read the argument through, and nothing here says that the Riemann hypothesis is
proved. This note treats one unit of the artifact: all \UnitInNote{} of its \UnitMembers{} declarations, Theorems 1
to~5 below. Their proofs are complete \emph{given} eleven results of other units, which are stated as Inputs~1
to~11 with the unit that owns each and are not proved here; among them are the three main theorems of the argument
(Inputs 1 and~2) and Landau's criterion (Input~10). Every statement of the unit and of the Inputs was compared with
its Lean text, argument by argument. The Lean text read is \PinText; Appendix~A is the concordance of \ConcDecls{}
declarations.}

\sect{The statements}{s:intro}
The objects in the five theorems are defined in §\ref{s:setting}: parameters and admissible data in
Definitions \ref{d:params} to~\ref{d:derived}, the remainders $\R{d}{M}{x}$ and $\rem{D}$ in
Definitions \ref{d:R} and~\ref{d:rem}, the three statements in Definition~\ref{d:statements}.

\statement{arith}{\phantomsection\label{t:1}\textbf{Theorem 1} (the main inequality).\enspace
\emph{Let $P$ be parameters and $D$ admissible data for~$P$, with number of lifts~$\Dm{D}$ and third scale~$\scb{D}$. Then}
\[
\Dm{D}\cdot\rem{D}\;\le\;(\log \scb{D})^2\cdot\bud{D}.
\]
\emph{Here $\rem{D}$ is the remainder of the weighted prime sums of~$D$, and
\m{\bud{D}=\fD{D}\cdot\aD{D}\cdot(\log \scb{D})^3} is its budget.}}
\leanline{\lean{star}}

\statement{arith}{\phantomsection\label{t:2}\textbf{Theorem 2} (the bound for the primes).\enspace
\emph{Let \m{P=(N_0,d)} be parameters and $\varepsilon$ a real number with \m{0<\varepsilon/6} and
\m{\varepsilon<\tfrac1{100}}. Then for every sufficiently large real number~$x$,}
\[
\R{d}{N_0}{x}\;\le\;x^{\,8\varepsilon+\frac12-1}.
\]}
\leanline{\lean{primeBound}}

\statement{arith}{\phantomsection\label{t:3}\textbf{Theorem 3} (the zero-free half-plane).\enspace
\emph{The half-plane statement holds: for every integer \m{M\ge1}, every Dirichlet character~$\chi$ modulo~$M$ and
every complex number~$s$ with \m{\re s>\tfrac12}, if $\chi$ is not the principal character modulo~$M$ or
\m{s\ne1}, then \m{L(s,\chi)\ne0}.}

\vspace{4pt}\phantomsection\label{t:4}\textbf{Theorem 4} (from the half-plane to the strip).\enspace
\emph{If the half-plane statement holds, then the strip statement holds: for every integer \m{M\ge1}, every
Dirichlet character~$\chi$ modulo~$M$ and every complex number~$s$ with \m{0<\re s}, \m{\re s<1} and
\m{L(s,\chi)=0}, one has \m{\re s=\tfrac12}.}

\vspace{4pt}\phantomsection\label{t:5}\textbf{Theorem 5} (from the half-plane to $\zeta$).\enspace
\emph{If the half-plane statement holds, then the Riemann hypothesis for~$\zeta$ holds: every complex number~$s$
with \m{\zeta(s)=0}, such that there is no natural number~$n$ with \m{s=-2(n+1)}, and with \m{s\ne1}, has
\m{\re s=\tfrac12}.}}
\leanline{\lean{grh}\sep \lean{zerosOnLine}\sep \lean{zetaZerosOnLine}}

Theorem~3 is the main theorem of the artifact. Theorems 3, 4 and~5 together give the five final theorems, which are
the subject of the note \emph{The five final theorems, and what one must trust}.

\subhead{The plan of the proof.} The argument attaches to the data~$D$ a finite-dimensional space of automorphic
forms, the \emph{rows}, and a linear map on it, the composite of $\Dm{D}$ theta lifts. Theorems A and~B of the
argument, taken together (Input~1), give linearly independent rows that almost span, have small index with respect
to the integral structures, and have short lifts: the squared covolume of the lifts is at most
\m{\exp\bigl(n\,(\hmain{D}-\Dm{D}\cdot\rem{D}+\text{error})\bigr)} at every conjugate, where $n$ is the number of the
rows and $\hmain{D}$ is the main amount. Theorem~C (Input~2) is a bound in the other direction, a consequence of the
product formula: integral forms are long. If the squared covolumes of the lifts of such rows are at most
\m{\exp(V)} at every conjugate, then $V$, plus twice a bound for the index, is at least
\m{\rowDim{D}\cdot(\hmain{D}-\text{error})}. The two bounds can hold together only if \m{\Dm{D}\cdot\rem{D}} is at most
the sum of the errors. This is Theorem~1, and it is all that the rest of the argument
needs from the automorphic side.

The remainder $\rem{D}$ is a combination of remainders of weighted prime sums at the first two scales of the
data (Input~6). For large~$x$ there are data with first scale \m{\scx{D}=x} (Input~7), and the second scale can be moved to a
point where the unwanted values are negligible (Input~9). Then Theorem~1, divided by~$\Dm{D}$ and by the other
positive factors, is Theorem~2. Landau's criterion (Input~10) gives Theorem~3, and the functional equations give
Theorems 4 and~5.

\sect{The setting}{s:setting}
\begin{conventions}\phantomsection\label{c:conventions}
The natural numbers are \m{0,1,2,\dots}; $\re s$ is the real part of a complex number~$s$; $\log$ is the natural
logarithm, and powers of positive real numbers with real exponents are defined through it; $\varphi$ is Euler's
function and $\Lambda$ von Mangoldt's function. A statement holds \kw{for every sufficiently large} real number
if there is a real number such that it holds for every real number at least as large. The Riemann zeta
function~$\zeta$, Dirichlet characters and their $L$-functions \m{L(s,\chi)} are those of Mathlib, the mathematical
library of Lean, at the commit that the artifact pins; the note \emph{The four statements: the Riemann hypothesis
and its generalisation to Dirichlet $L$-functions} says what determines them. A modulus is an integer \m{M\ge1}. The letter $\nu$ stands for an integer, and for a natural number where a weight is defined; it is never a character of the ideles (Convention~0.6). A bare $x$ is a free real variable (Convention~0.3), and the scales of the data $D$ are written $x(D)$, $x_1(D)$, $b(D)$.
Parameters are called~$P$ and admissible data for~$P$ are called~$D$; a statement about~$D$ is a statement for
every~$P$ and all such~$D$, and the components of $P$ and of~$D$ are named where they are used.
\end{conventions}

\subhead{The three statements.}
\begin{definition}[three propositions]\label{d:statements}
\textup{(1)} The \kw{Riemann hypothesis for $\zeta$}: for every complex number~$s$, if \m{\zeta(s)=0}, if there
is no natural number~$n$ with \m{s=-2(n+1)}, and if \m{s\ne1}, then \m{\re s=\tfrac12}.
\textup{(2)} The \kw{half-plane statement}: for every integer \m{M\ge1}, every Dirichlet character~$\chi$ modulo~$M$
and every complex number~$s$ with \m{\re s>\tfrac12}, if $\chi$ is not the principal character modulo~$M$ or
\m{s\ne1}, then \m{L(s,\chi)\ne0}.
\textup{(3)} The \kw{strip statement}: for every integer \m{M\ge1}, every Dirichlet character~$\chi$ modulo~$M$ and
every complex number~$s$, if \m{0<\re s}, \m{\re s<1} and \m{L(s,\chi)=0}, then \m{\re s=\tfrac12}.
\end{definition}
\leanline{\lean{rh}\sep \lean{halfPlane}\sep \lean{strip}\sep another unit; see the note \emph{The four statements}}

\subhead{Parameters and data.}
\begin{definition}[parameters]\label{d:params}
\kw{Parameters} are a pair \m{P=(N_0,d)} of natural numbers such that $18480$ divides~$N_0$, $d$ is a prime number,
$N_0$ divides \m{2(d+1)}, and $N_0$ does not divide \m{d+1}. For parameters \m{P=(N_0,d)} we write
\m{K=\QQ(\sqrt{-d})}, a subfield of~$\CC$, and $\chi_K$ for the quadratic character of the field \m{\ZZ/d}, a
Dirichlet character modulo~$d$; it is the Legendre symbol \m{(\,\cdot\,/d)} when $d$ is odd, as it is for all
parameters (Lemma~\ref{l:params}).
\end{definition}
\leanline{\lean{params}\sep unit 152}

The last two conditions say that \m{d\equiv N_0/2-1} modulo~$N_0$; they force \m{N_0\ge1}. Since
\m{18480=2^4\cdot3\cdot5\cdot7\cdot11}, one has \m{d\equiv7} modulo~$8$ (Lemma~\ref{l:params}); so $K$ is an
imaginary quadratic field in which $2$ splits, and $\chi_K$ is its quadratic character.

\begin{definition}[admissible data]\label{d:data}
Let \m{P=(N_0,d)} be parameters. The artifact attaches to~$P$ a real number~$\cP{P}$, of which this note uses only
Input~5. \kw{Admissible data for~$P$} are a record~$D$ with the following components.
\begin{enumerate}[label=\textup{(\alph*)},leftmargin=2.3em,itemsep=2pt]
\item Three real numbers \m{\log \scx{D}}, \m{\log \scxo{D}} and \m{\log \scb{D}}.
Their exponentials $\scx{D}$, $\scxo{D}$ and~$\scb{D}$ are the
first, second and third \kw{scale} of~$D$.
\item Two natural numbers: $\Dm{D}$, the \kw{number of lifts}, and $\Dr{D}$, a rank. By condition~(ii) below, $\Dm{D}$ is a power
of~$2$; put
\[
\None{D}=N_0\cdot \Dm{D}^{22}\cdot3^{\lfloor\log_3(22\log_2\Dm{D})\rfloor+5},
\]
and let $\ED{D}$ be the field generated over~$K$ by the roots of unity of order~$\None{D}$.
\item A finite set $\thick{D}$ of infinite places of the field~$\ED{D}$, the \kw{thick} places.
\item Two prime numbers and an element of the ring of integers of~$K$.
\end{enumerate}
They are subject to the following conditions.
\begin{enumerate}[label=\textup{(\roman*)},leftmargin=2.3em,itemsep=2pt]
\item \emph{The scales.}
\[
\begin{gathered}
10^{12}\,\cP{P}\le\log\log \scb{D},\qquad 100\,\log\log \scb{D}\le\log \scx{D},\\
10^{6}\,\log \scx{D}\le\log \scxo{D},\qquad 10^{6}\,\log \scxo{D}\le\log \scb{D},\\
6\,\log\log \scxo{D}\le\log\log \scb{D}.
\end{gathered}
\]
\item \emph{The number of lifts and the rank.} $\Dm{D}$ is a power of~$2$ with exponent at least~$1$,
\m{\sqrt{\scx{D}}\le \Dm{D}\le2\sqrt{\scx{D}}} and \m{22\log_2\Dm{D}\le32\log \scx{D}}; $\Dr{D}$ is a power of~$3$ with exponent at least~$1$, and
\m{\scx{D}^{5}\le \Dr{D}\le3\,\scx{D}^{5}}.
\item \emph{Further conditions,} on the classes of the thick places (a congruence, and the absence of certain
relations between three of them), on the two primes and the integer of~(d) (the primes are distinct, satisfy
congruences and lie between $\scb{D}^{10}$ and~$\scb{D}^{11}$, and the norm of the integer is their product), and the existence
of primes in certain short intervals.
\end{enumerate}
The conditions~(iii) are stated in full in the note on the data of the proof. In this note they matter only through
the Inputs, and through one fact that was checked in the Lean text: \m{\log \scxo{D}} occurs in the conditions~(i) and
nowhere else.
\end{definition}
\leanline{\lean{adm}\sep \lean{cP}\sep \lean{None}\sep \lean{ED}\sep unit 152}

The conditions~(i) put the scales far apart, \m{(\log \scb{D})^{100}\le \scx{D}}, \m{\scx{D}^{10^6}\le \scxo{D}}, \m{\scxo{D}^{10^6}\le \scb{D}},
and make all of them very large (Lemma~\ref{l:sizes}). The number $\None{D}$ and the field $\ED{D}$ depend only
on $P$ and~$\Dm{D}$. The formal record has one more component, a natural number called the exponent, which a condition
makes equal to \m{22\log_2\Dm{D}}; this note does not give it a letter (§\ref{s:notes}, note~3). In the formal text the
quotient \m{\None{D}/N_0} is the degree of a tower of fields, a function of the exponent.

\begin{definition}[quantities derived from the data]\label{d:derived}
Let \m{P=(N_0,d)} be parameters and $D$ admissible data for~$P$, with scales $\scx{D}$, $\scxo{D}$,~$\scb{D}$.
\begin{enumerate}[label=\textup{(\alph*)},leftmargin=2.3em,itemsep=2pt]
\item Every infinite place of~$\ED{D}$ is the place of exactly one embedding \m{\ED{D}\to\CC} that is the identity
on~$K$. The Galois group of $\ED{D}$ over~$K$ permutes these embeddings simply transitively, and its action on the
roots of unity identifies it with \m{(\ZZ/\None{D})^{\times}}. (Restriction to the field generated over~$\QQ$ by these roots of unity is an isomorphism onto the Galois group
of that field over~$\QQ$, because the quadratic field~$K$ is not contained in it: the prime~$d$ ramifies in~$K$,
and it does not ramify in a field of roots of unity of order prime to~$d$. And $d$ is prime to~$\None{D}$: it is
prime to~$N_0$ by Lemma~\ref{l:params}, and $\None{D}$ has the same prime divisors as~$N_0$, since $\Dm{D}$ is a power
of~$2$ and $6$ divides~$N_0$.) The artifact chooses one of the embeddings once.
Each infinite place~$\tau$ of~$\ED{D}$ then has a \kw{class} \m{\cl{D}{\tau}\in(\ZZ/\None{D})^{\times}}, the chosen
embedding has class~$1$, and \m{\tau\mapsto\cl{D}{\tau}} is a bijection. For an integer~$\nu$ prime to~$\None{D}$,
\m{\nu\cdot\tau} is the place whose class is $\nu$ times that of~$\tau$; for the other integers~$\nu$ we put
\m{\nu\cdot\tau=\tau}.
\item $\tops{D}$, the set of \kw{tops}, is the set of the infinite places~$\tau$ of~$\ED{D}$ whose class
\m{\cl{D}{\tau}} is congruent to~$1$ modulo~$N_0$.
\item $\fD{D}$ is the number of infinite places of the field~$\ED{D}$.
\item $\aD{D}$ is the least multiple of~$\None{D}$ that is at least \m{\scx{D}\,\scb{D}}; it is a natural number.
\item \m{\bud{D}=\fD{D}\cdot\aD{D}\cdot(\log \scb{D})^{3}} is the \kw{budget}.
\end{enumerate}
\end{definition}
\leanline{\lean{tops}\sep \lean{fD}\sep \lean{aD}\sep \lean{budget}\sep unit 152}

\subhead{The weighted prime sums.}
\begin{definition}[the weight, and the remainder of the prime sum]\label{d:R}
Let $d$ be a prime number and \m{M\ge1} an integer; \m{K=\QQ(\sqrt{-d})}, and $\chi_K$ is the quadratic
character of the field \m{\ZZ/d}, a Dirichlet character modulo~$d$ (the Legendre symbol \m{(\,\cdot\,/d)} when $d$
is odd).
\begin{enumerate}[label=\textup{(\alph*)},leftmargin=2.3em,itemsep=2pt]
\item For a natural number~$\nu$ put \m{\lam{d}{M}{\nu}=\Lambda(\nu)\bigl(1-\chi_K(\nu)\bigr)} if \m{\nu\equiv1}
modulo~$M$, and \m{\lam{d}{M}{\nu}=0} otherwise.
\item For a Dirichlet character~$\chi$ modulo an integer \m{M\ge1}, let \m{s\mapsto\entireLs{\chi}} be the
following entire function: \m{\entireLs{\chi}=L(s,\chi)} if $\chi$ is not principal; if $\chi$ is principal,
\m{\entireLs{\chi}=(s-1)L(s,\chi)} for \m{s\ne1}, and \m{\entireLone{\chi}=\prod_{p\mid M}(1-p^{-1})}, the
product over the primes~$p$ dividing~$M$.
\item Put
\[
\kap{d}{M}=\frac1{\varphi(M)}\,\re\Bigl(\sum_{\xi}\Bigl(\frac{\tilde L'}{\tilde L}(1,\xi\chi_K)
-\frac{\tilde L'}{\tilde L}(1,\xi)\Bigr)-1\Bigr),
\]
where $\xi$ runs over the Dirichlet characters modulo~$M$, \m{\xi\chi_K} is the character modulo~$Md$ with the
values \m{\xi(\nu)\chi_K(\nu)}, and \m{\frac{\tilde L'}{\tilde L}(s,\chi)} is the logarithmic derivative of
\m{s\mapsto\entireLs{\chi}}.
\item For a real number \m{x>0} put
\[
\R{d}{M}{x}=\sum_{1\le\nu\le x}\lam{d}{M}{\nu}\Bigl(\frac1\nu-\frac1x\Bigr)-\frac{\log x}{\varphi(M)}-\kap{d}{M}.
\]
\end{enumerate}
\end{definition}
\leanline{\lean{lam}\sep \lean{R}\sep unit 151\sep \lean{entireL}\sep unit 106, background}

The weight $\lam{d}{M}{\nu}$ is never negative, and it vanishes unless $\nu$ is a power of a prime~$p$ with
\m{\chi_K(p)\ne1}; when \m{d\equiv3} modulo~$4$, as for every prime~$d$ to which the theorems below are applied, these
are the primes that do not split in~$K$. The sum in~(d) is a smoothed count of such prime powers in the class of~$1$ modulo~$M$; the term
\m{\log x/\varphi(M)+\kap{d}{M}} is its expected value, and $\R{d}{M}{x}$ is the \kw{remainder}. That it tends
to~$0$ is a form of the prime number theorem in progressions; under the generalised Riemann hypothesis it is
bounded by a constant times \m{x^{-1/2}}, by the explicit formula, and Theorem~2 is an upper bound of nearly this
strength.

\begin{definition}[the remainder of the data]\label{d:rem}
Let \m{P=(N_0,d)} be parameters and $D$ admissible data for~$P$, with scales $\scx{D}$, $\scxo{D}$,~$\scb{D}$; in this definition
$\tau$ is an infinite place of~$\ED{D}$.
\begin{enumerate}[label=\textup{(\alph*)},leftmargin=2.3em,itemsep=2pt]
\item The \kw{nominal weight} \m{\nomW{D}{\tau}} of~$\tau$ is \m{\aD{D}/\scx{D}} if \m{\tau\in\thick{D}}, and it is
\m{\aD{D}/\scxo{D}} otherwise.
\item For a real number~$z$,
\[
\Ask{D}{z}{\tau}=\sum_{\nu=1}^{\aD{D}}\lam{d}{1}{\nu}\,\Bigl(\frac{z}{\nu}-\nomW{D}{(-\nu)\cdot\tau}\Bigr)_{\!+},
\]
where \m{(\,\cdot\,)_{+}} is the positive part of a real number.
\item For the prime number~$d$, \m{\chowla{d}} is Euler's constant minus the real part of the logarithmic
derivative of \m{s\mapsto\entireLs{\chi_K}} at \m{s=1}. Here $\chi_K$ is not principal, so this is the real part of
\m{(L'/L)(1,\chi_K)}.
\item \m{\mainAsk{D}=\aD{D}\cdot\bigl(\log\aD{D}-1-\chowla{d}-\fD{D}^{-1}\sum_{\tau'}\log\nomW{D}{\tau'}\bigr)},
where $\tau'$ runs over all infinite places of~$\ED{D}$; it is the \kw{expected value} of the sums in~(b) for
\m{z=\aD{D}}.
\item The \kw{remainder} of~$D$ is
\[
\rem{D}=\sum_{\tau\in\tops{D}}\bigl(\Ask{D}{\aD{D}}{\tau}-\mainAsk{D}\bigr).
\]
\end{enumerate}
\end{definition}
\leanline{\lean{nomW}\sep \lean{Ask}\sep unit 152\sep \lean{mainAsk}\sep \lean{rem}\sep unit 151\sep
\lean{chowla}\sep unit 106, background}

So $\rem{D}$ compares weighted prime sums at the tops with their expected value. A lemma of unit~151 shows that
\m{-1-\chowla{d}=\kap{d}{1}}; the overview of this handout writes the expected value with~$\kap{d}{1}$. Input~6
expresses $\rem{D}$ through the remainders of Definition~\ref{d:R}, and nothing else about
Definition~\ref{d:rem} is used in this note.

\subhead{Notions of later chapters.} Theorem~1 mentions none of the following notions, but its proof passes
through them. They are named here with all their arguments; their constructions belong to later chapters and are
not needed, because the proof only uses that the conclusion of Input~1 fits the hypothesis of Input~2.

\begin{definition}[rows, lifts, covolume, index, main amount]\label{d:later}
Let $P$ be parameters and $D$ admissible data for~$P$, with number of lifts~$\Dm{D}$ and rank~$\Dr{D}$.
\begin{enumerate}[label=\textup{(\alph*)},leftmargin=2.3em,itemsep=2pt]
\item $\rows{D}$ is a complex vector space, the space of \kw{rows}: certain cusp forms on a unitary group of
rank~$\Dr{D}$ attached to the field~$\ED{D}$. Its dimension over~$\CC$ is \m{\rowDim{D}}; in the library the dimension
of a space that is not finite-dimensional is~$0$.
\item The data have \kw{$\Dm{D}$-fold twists}~$\vartheta$: certain characters, defined in later chapters, each with a
finite set of primes of~$\ED{D}$ at which it ramifies. Among them is the trivial twist~$1$, which ramifies at no
prime. For an $\Dm{D}$-fold twist~$\vartheta$, \m{\Lift{D}{\vartheta}} is a $\CC$-linear map from $\rows{D}$ to a space
of algebraic modular forms: the composite of $\Dm{D}$ theta lifts, formed with the first $\Dm{D}$ powers of~$\vartheta$,
restricted to the rows.
\item A space of algebraic modular forms of the kind in~(b) is given by a CM field, a rank, a set of level places,
a level subgroup with a character of it, and weights at the infinite places. For forms \m{F_1,\dots,F_n} in one
such space and an automorphism~$\sigma$ of the field~$\CC$, the \kw{squared covolume}
\m{\covol{F_1,\dots,F_n}{\sigma}} is the absolute value of the determinant of the matrix of the pairings of the
$\sigma$-conjugates of \m{F_1,\dots,F_n}; it is a real number \m{\ge0}.
\item For a natural number~$n$, rows \m{y_1,\dots,y_n} in~$\rows{D}$ and a real number~$\delta$, the relation
\m{\indexLe{D}{y_1,\dots,y_n}{\delta}} means: there are real numbers~$\eta_\ell$, one for each natural
number~$\ell$, with \m{\sum_{\ell\in J}\eta_\ell\le\delta} for every finite set~$J$ of natural numbers, such that
for every prime number~$\ell$, every $\ell$-adic place of~$\CC$ and every $\Dm{D}$-fold twist~$\vartheta$ to which
$\ell$ is prime (that is, $\ell$ is prime to the norms of the primes at which $\vartheta$ ramifies), some
\m{n\times n} matrix over the valuation ring of that place, with determinant of absolute value at least
\m{\exp(-\eta_\ell)}, transforms \m{y_1,\dots,y_n} into rows all of whose images under \m{\Lift{D}{\vartheta}} are
integral at that place.
\item $\hmain{D}$ is a real number, the \kw{main amount}: the product of $\Dm{D}$, of the rank \m{\Dr{D}+1-\Dm{D}}, and of a
real number called the cost of one direction, which is defined from the weights and the prime sums of~$D$.
\end{enumerate}
\end{definition}
\leanline{\lean{rows}\sep \lean{rowDim}\sep unit 140\sep \lean{Lift}\sep unit 143\sep \lean{covol}\sep unit 102,
background\sep \lean{indexLe}\sep \lean{hmain}\sep unit 151}

\sect{The main inequality}{s:star}
The proof compares two bounds for one quantity. For rows \m{y_1,\dots,y_n} as in Input~1 below, let $V$ be the
logarithm of a bound for the squared covolumes of their lifts at all automorphisms of~$\CC$, and $\delta$ a bound
for their index. Input~1 bounds \m{V+2\delta} from above, Input~2 bounds it from below, and up to errors that are
multiples of the budget,
\[
\hmain{D}\;\le\;\frac{V+2\delta}{\rowDim{D}}\;\le\;\hmain{D}-\Dm{D}\cdot\rem{D}.
\]
The inequality between the two ends is Theorem~1. This display is a guide only; it stands in a comment at the head
of the file of the unit, in which the authors also list the inputs and the sizes of the errors. The proof follows
the Inputs.

\statement{inksoft}{\textbf{Inputs from other units.} The following results are taken for granted. They are stated
as the dossier of this unit reads them, and each was compared with its Lean text.

\vspace{4pt}\phantomsection\label{i:1}\textbf{Input 1} (Theorems A and B together; unit 151).\enspace
\emph{Let $P$ be parameters and $D$ admissible data for~$P$, with number of lifts~$\Dm{D}$ and third scale~$\scb{D}$. Then
there are a natural number~$n$ and rows \m{y_1,\dots,y_n} in~$\rows{D}$ such that: \m{y_1,\dots,y_n} are linearly
independent over~$\CC$; \m{(1-4/\scb{D}^{2})\cdot\rowDim{D}\le n}; there is a real number~$\delta$ with
\m{\delta\le5\cdot\rowDim{D}\cdot\log(1+\scb{D}^{2})\cdot\bud{D}} such that \m{\indexLe{D}{y_1,\dots,y_n}{\delta}}; and
for every automorphism~$\sigma$ of~$\CC$,}
\[
\begin{gathered}
\covol{\Lift{D}{1}y_1,\dots,\Lift{D}{1}y_n}{\sigma}\\
\le\;\exp\bigl(n\cdot(\hmain{D}-\Dm{D}\cdot\rem{D}+\log \scb{D}\cdot\bud{D})\bigr).
\end{gathered}
\]

\vspace{4pt}\phantomsection\label{i:2}\textbf{Input 2} (Theorem C: integral forms are long; unit 151).\enspace
\emph{Let $P$ be parameters, $D$ admissible data for~$P$ with third scale~$\scb{D}$, $n$ a natural number,
\m{y_1,\dots,y_n} rows in~$\rows{D}$, and $\delta$ and~$V$ real numbers. Assume that \m{y_1,\dots,y_n} are linearly
independent over~$\CC$, that \m{(1-4/\scb{D}^{2})\cdot\rowDim{D}\le n}, that \m{\indexLe{D}{y_1,\dots,y_n}{\delta}}, and
that \m{\covol{\Lift{D}{1}y_1,\dots,\Lift{D}{1}y_n}{\sigma}\le\exp(V)} for every automorphism~$\sigma$ of~$\CC$.
Then}
\[
\begin{gathered}\rowDim{D}\cdot\bigl(\hmain{D}-3\cdot\log \scb{D}\cdot\bud{D}\bigr)\\\le\;V+2\delta.\end{gathered}
\]}
\leanline{\lean{familyAB}\sep \lean{longC}}

In the words of the argument, Input~1 puts together Theorem~A (short forms: there are integral forms whose lifts
are short, with a bound for the index) and Theorem~B (enough forms: few translates of them almost span the rows).
Input~2 is Theorem~C. Three smaller inputs are needed as well.

\statement{inksoft}{\phantomsection\label{i:3}\textbf{Input 3} (a crude bound; unit 151).\enspace
\emph{Let $P$ be parameters and $D$ admissible data for~$P$, with number of lifts~$\Dm{D}$ and scales $\scx{D}$ and~$\scb{D}$. Then
\m{|\hmain{D}|+\Dm{D}\cdot|\rem{D}|\le \scb{D}^{2}\cdot \scx{D}^{99}}.}

\vspace{4pt}\phantomsection\label{i:4}\textbf{Input 4} (there are rows; unit 148).\enspace
\emph{Let $P$ be parameters and $D$ admissible data for~$P$. Then \m{\rowDim{D}\ne0}; in particular $\rows{D}$ has
finite dimension.}

\vspace{4pt}\phantomsection\label{i:5}\textbf{Input 5} (the constant and the places; unit 152).\enspace
\emph{Let \m{P=(N_0,d)} be parameters. Then \m{64\le\cP{P}}, \m{N_0\le\cP{P}}, \m{d\le\cP{P}} and
\m{\varphi(N_0)\le\cP{P}}. Let moreover $D$ be admissible data for~$P$. Then $N_0$ divides~$\None{D}$,}
\[
\fD{D}=\frac{\None{D}}{N_0}\cdot\varphi(N_0)\qquad\text{\emph{and}}\qquad\#\tops{D}=\frac{\None{D}}{N_0}.
\]
\emph{The formal text states the two equations for every exponent of the tower; they are used here for the
exponent of~$D$.}}
\leanline{\lean{crude}\sep \lean{rowsNe}\sep \lean{cGe}\sep \lean{cardPlaces}\sep \lean{cardTops}}

\begin{lemma}[sizes]\label{l:sizes}
Let \m{P=(N_0,d)} be parameters and $D$ admissible data for~$P$, with scales $\scx{D}$, $\scxo{D}$,~$\scb{D}$. Then
\begin{enumerate}[label=\textup{(\alph*)},leftmargin=2.3em,itemsep=2pt]
\item $\log \scb{D}\ge10^{6}\log \scxo{D}\ge10^{12}\log \scx{D}$ and $\log \scx{D}\ge100\log\log \scb{D}\ge10^{14}\,\cP{P}$, and also
$\cP{P}\ge64$; in particular $\scx{D}\ge2$, $\log \scb{D}\ge35$, $\scx{D}^{99}\le \scb{D}$ and $\log \scb{D}\le \scx{D}^{1/100}$;
\item $\aD{D}\ge \scx{D}\,\scb{D}$, $\fD{D}\ge1$, and $\bud{D}\ge \scb{D}\,(\log \scb{D})^{3}>0$;
\item $4\cdot\log \scb{D}\cdot\bud{D}+4\,\scx{D}^{99}+10\cdot\log(1+\scb{D}^{2})\cdot\bud{D}\;\le\;(\log \scb{D})^{2}\cdot\bud{D}$.
\end{enumerate}
\end{lemma}
\begin{proof}
(a) The four inequalities between the logarithms are conditions~(i) of Definition~\ref{d:data}, and
$\cP{P}\ge64$ is in Input~5. Hence $\log \scx{D}\ge64\cdot10^{14}$, which gives $\scx{D}\ge2$; then
$\log \scb{D}\ge10^{12}\log \scx{D}\ge35$; $99\log \scx{D}\le10^{12}\log \scx{D}\le\log \scb{D}$ gives $\scx{D}^{99}\le \scb{D}$; and
$100\log\log \scb{D}\le\log \scx{D}$ gives $\log \scb{D}\le \scx{D}^{1/100}$.

(b) The first inequality is the definition of~$\aD{D}$. A number field has an infinite place, so $\fD{D}\ge1$.
With $\scx{D}\ge1$ this gives $\bud{D}=\fD{D}\cdot\aD{D}\cdot(\log \scb{D})^{3}\ge \scb{D}\,(\log \scb{D})^{3}$, which is positive
by~(a).

(c) By (a) and~(\scb{D}), $4\scx{D}^{99}\le4\scb{D}\le \scb{D}\,(\log \scb{D})^{3}\le\bud{D}$. As \m{\scb{D}\ge1} and \m{\log \scb{D}\ge\log2}, also
$\log(1+\scb{D}^{2})\le\log(2\scb{D}^{2})\le3\log \scb{D}$. So the left side of~(c) is at most
\[
(4\log \scb{D}+1+30\log \scb{D})\cdot\bud{D}\;\le\;35\,\log \scb{D}\cdot\bud{D},
\]
and \m{35\le\log \scb{D}}.
\end{proof}
\leanline{in the formal text these comparisons are made by a calculus of sizes of unit 152, \lean{sizeLe}; see
§\ref{s:notes}, note~5}

\begin{proof}[Proof of Theorem 1]
Let $n$, \m{y_1,\dots,y_n} and~$\delta$ be as in Input~1, and put
\[
\hA{}=\hmain{D}-\Dm{D}\cdot\rem{D}+\log \scb{D}\cdot\bud{D}.
\]
The last property in Input~1 says that the squared covolumes of \m{\Lift{D}{1}y_1,\dots,\Lift{D}{1}y_n} are at most
\m{\exp(n\,\hA{})}. So Input~2 applies to $n$, \m{y_1,\dots,y_n}, $\delta$ and \m{V=n\,\hA{}}, and gives
\begin{equation}\label{e:both}
\begin{aligned}
&\rowDim{D}\cdot\bigl(\hmain{D}-3\cdot\log \scb{D}\cdot\bud{D}\bigr)\;\le\;n\,\hA{}+2\delta\\
&\qquad=\;\rowDim{D}\cdot\hA{}+\bigl(\rowDim{D}-n\bigr)\cdot(-\hA{})+2\delta .
\end{aligned}
\end{equation}
The two bounds are compared on families of different sizes, $n$ and $\rowDim{D}$, and the middle term on the
right pays for the difference. It is small for the following reasons. By Input~4 the space of rows has finite
dimension, so the linearly independent rows \m{y_1,\dots,y_n} are at most $\rowDim{D}$ in number, and with the
second property in Input~1,
\[
0\;\le\;\rowDim{D}-n\;\le\;\frac{4}{\scb{D}^{2}}\cdot\rowDim{D}.
\]
By Lemma~\ref{l:sizes}(b), \m{\log \scb{D}\cdot\bud{D}>0}. With Input~3 and \m{\Dm{D}\ge0},
\[
\begin{aligned}
-\hA{}&=-\hmain{D}+\Dm{D}\cdot\rem{D}-\log \scb{D}\cdot\bud{D}\\
&\le\;|\hmain{D}|+\Dm{D}\cdot|\rem{D}|\;\le\;\scb{D}^{2}\,\scx{D}^{99}.
\end{aligned}
\]
Hence \m{(\rowDim{D}-n)\cdot(-\hA{})\le4\,\scx{D}^{99}\cdot\rowDim{D}}. The bound for~$\delta$ in Input~1 gives
\m{2\delta\le10\cdot\log(1+\scb{D}^{2})\cdot\bud{D}\cdot\rowDim{D}}. Putting these into~\eqref{e:both},
\[
\hmain{D}-3\cdot\log \scb{D}\cdot\bud{D}\;\le\;\hA{}+4\,\scx{D}^{99}+10\cdot\log(1+\scb{D}^{2})\cdot\bud{D},
\]
after division by \m{\rowDim{D}}, which is positive by Input~4.
Inserting the definition of~$\hA{}$, the main amount $\hmain{D}$ cancels, and
\[
\begin{aligned}
\Dm{D}\cdot\rem{D}\;&\le\;4\cdot\log \scb{D}\cdot\bud{D}+4\,\scx{D}^{99}\\
&\qquad{}+10\cdot\log(1+\scb{D}^{2})\cdot\bud{D}\\
&\le\;(\log \scb{D})^{2}\cdot\bud{D},
\end{aligned}
\]
the last step by Lemma~\ref{l:sizes}(c).
\end{proof}
\leanline{\lean{star}}

The three terms of the error have separate sources: \m{4\cdot\log \scb{D}\cdot\bud{D}} is the error of length in Input~1
and the error of covolume in Input~2, \m{4\scx{D}^{99}} pays for the dimensions that the rows do not span, and the last
term is the index. All three are far below \m{(\log \scb{D})^{2}\cdot\bud{D}}; the exponents $99$ and~$2$ are slack.

\sect{The bound for the primes}{s:primes}
\begin{definition}[polar functions]\label{d:polar}
Let \m{\cE\colon\RR\to\RR} and \m{\omega\colon\CC\to\CC} be functions. We say that
\polarWords{\cE}{\omega}, and write \m{\polar{\cE}{\omega}}, if $\cE$ is measurable and there is a function
\m{\Phi\colon\CC\to\CC}, meromorphic on~$\CC$ and defined at every point (at a pole its value is arbitrary), such
that:
the function \m{v\mapsto|\cE{}(v)|\exp(-v)} is bounded on \m{v\ge0};
\m{\Phi(s)=\int_0^{\infty}\cE{}(v)\exp((1-s)v)\,dv} for \m{\re s>2};
the function \m{s\mapsto(s-1)\Phi(s)}, whose value at \m{s=1} is~$0$ whatever $\Phi(1)$ is, is analytic at every
point of the closed half-plane \m{\re s\ge1}; and for every complex number~$\rho$ with \m{\re\rho>0},
\m{(s-\rho)\Phi(s)\to\omega(\rho)} as \m{s\to\rho}, \m{s\ne\rho}.
\end{definition}
\leanline{\lean{isPolar}\sep unit 107, background}

So $\Phi$ continues the Laplace transform of~$\cE$. The third condition is to be read with its clause on the value
at~$1$: it says that $\Phi$ is analytic at every point \m{s\ne1} of \m{\re s\ge1}, and that near \m{s=1} it is the
quotient by \m{s-1} of an analytic function that vanishes at~$1$. So the singularity at~$1$ is removable, $\Phi$
has no pole on \m{\re s\ge1}, and \m{\omega(\rho)=0} for \m{\re\rho\ge1}. The clause cannot be dropped. Without
it the constant function~$1$, whose transform is \m{1/(s-1)}, would be polar, and Input~9 below would be false
for the constant functions $1$ and~$0$: it would give \m{1\le\exp(-v_0)} for large~$v_0$. In the half-plane
\m{\re s>0} the poles of $\Phi$ are at most simple, with residue~$\omega(\rho)$ at~$\rho$. The definition, with a
remark that derives these consequences, is also in Chapter~II of this handout, the note \emph{From a one-sided
bound for primes to a zero-free half-plane}, in the same letters.

\statement{inksoft}{\textbf{Inputs from other units.}

\vspace{4pt}\phantomsection\label{i:6}\textbf{Input 6} (the remainder of the data, through the remainders of the prime sums; unit 151).\enspace
\emph{Let \m{P=(N_0,d)} be parameters and $D$ admissible data for~$P$, with first scale~$\scx{D}$ and second
scale~$\scxo{D}$. Then}
\[
\begin{aligned}
\rem{D}=\aD{D}\cdot\Bigl(&\#\thick{D}\cdot\bigl(\R{d}{N_0}{\scx{D}}-\R{d}{N_0}{\scxo{D}}\bigr)\\
&+\#\tops{D}\cdot\R{d}{1}{\scxo{D}}\Bigr).
\end{aligned}
\]

\vspace{4pt}\phantomsection\label{i:7}\textbf{Input 7} (data exist; unit 152).\enspace
\emph{Let \m{P=(N_0,d)} be parameters and $\varepsilon$ a real number with \m{0<\varepsilon/6} and
\m{\varepsilon<\tfrac1{100}}. Then for every sufficiently large real number~$x$ there are admissible data~$D$
for~$P$ with first scale \m{\scx{D}=x}, whose second scale~$\scxo{D}$ and third scale~$\scb{D}$ satisfy}
\[
\begin{gathered}
\log \scxo{D}=x^{\varepsilon/6},\\ (10^{6}\log x)^{2}\le\log \scxo{D},\\
\log\log \scb{D}=\varepsilon\log x,\\ \frac{\None{D}/N_0}{\log \scb{D}}\le\#\thick{D}.
\end{gathered}
\]}
\leanline{\lean{remEq}\sep \lean{existsData}}

Input~6 is where the prime sums of Definition~\ref{d:rem} become the remainders of Definition~\ref{d:R}: as
$\tau$ runs over the tops, the places \m{(-\nu)\cdot\tau} run over the classes \m{\equiv-\nu} modulo~$N_0$, which
contain the thick places exactly when \m{\nu\equiv1} modulo~$N_0$. So the thick places see the weight
\m{\lam{d}{N_0}{\nu}} at the first scale, every place sees \m{\lam{d}{1}{\nu}} at the second scale, and the main
terms are the expected values. Input~7 provides the data, and the next two inputs make it possible to move the second scale.

\statement{inksoft}{\phantomsection\label{i:8}\textbf{Input 8} (the remainders are polar; unit 151).\enspace
\emph{Let $d$ be a prime number and \m{M\ge1} an integer such that \m{d\equiv7} modulo~$8$ and $M$ is prime
to~$d$. For a complex number \m{\rho\ne0,1} let}
\[
\res{d}{M}{\rho}=\frac{1}{\varphi(M)\,\rho\,(1-\rho)}\sum_{\xi}\bigl(\ord_{\rho}\entireL{\xi}
-\ord_{\rho}\entireL{\xi\chi_K}\bigr),
\]
\emph{where $\xi$ runs over the Dirichlet characters modulo~$M$, \m{\xi\chi_K} is as in Definition~\ref{d:R}, and
$\ord_{\rho}$ is the order of vanishing at~$\rho$ of an entire function that is not zero; and let
\m{\res{d}{M}{0}=\res{d}{M}{1}=0}. Then, as a function of the real variable~$v$,
\polarWords{v\mapsto\R{d}{M}{\exp(v)}}{\rho\mapsto\res{d}{M}{\rho}}.}

\vspace{4pt}\phantomsection\label{i:9}\textbf{Input 9} (a good point in every long window; unit 151).\enspace
\emph{Let $\cE_1$, $\cE_2$ be functions from $\RR$ to~$\RR$ and $\omega_1$, $\omega_2$ functions from $\CC$
to~$\CC$ such that \polarWords{\cE_1}{\omega_1} and \polarWords{\cE_2}{\omega_2}. Then for every sufficiently
large real number~$v_0$ and every real number~$X$ with \m{0\le X\le\exp(v_0)} there is a real number~$v_1$ with
\m{v_0\le v_1\le v_0^{2}} and}
\[
\cE_1(v_1)-X\cdot\cE_2(v_1)\;\le\;\exp(-v_0).
\]}
\leanline{\lean{polarR}\sep \lean{window}}

In Input~8 the formal text gives the residue function by the one formula, for every complex~$\rho$; a quotient
by~$0$ is~$0$ in Lean, which is the value at \m{\rho=0} and \m{\rho=1} stated above. The values of the residue
function play no role in this note, which uses only that the remainders are polar.

\begin{lemma}[the parameters]\label{l:params}
Let \m{P=(N_0,d)} be parameters. Then \m{d\equiv7} modulo~$8$, and $N_0$ is prime to~$d$.
\end{lemma}
\begin{proof}
As $16$ divides $18480$, which divides~$N_0$, and $N_0$ divides \m{2(d+1)}, the number~$8$ divides \m{d+1}. Let
$p$ be a prime that divides~$N_0$. Then $p$ divides \m{2(d+1)}. If \m{p=d}, then $d$ divides~$2$, which is
impossible since \m{d\equiv7} modulo~$8$. Hence no prime divides both $N_0$ and~$d$.
\end{proof}
\leanline{\lean{heeg}\sep \lean{coprime}\sep unit 152; the proof here is direct}

\begin{lemma}[moving the second scale]\label{l:move}
Let $P$ be parameters, $D$ admissible data for~$P$ with scales $\scx{D}$, $\scxo{D}$,~$\scb{D}$, and $v_1$ a real number with
\m{10^{6}\log \scx{D}\le v_1\le\log \scxo{D}}. Let $D'$ be the record whose second scale is \m{\exp(v_1)} and whose other
components are those of~$D$. Then $D'$ are admissible data for~$P$. Moreover \m{\None{D'}=\None{D}},
\m{\ED{D'}=\ED{D}}, \m{\thick{D'}=\thick{D}}, \m{\tops{D'}=\tops{D}}, \m{\fD{D'}=\fD{D}}, \m{\aD{D'}=\aD{D}} and
\m{\bud{D'}=\bud{D}}.
\end{lemma}
\begin{proof}
Only the conditions~(i) of Definition~\ref{d:data} mention $\log \scxo{D}$, and of these the last three. For~$D'$
the first of the three is the lower bound for~$v_1$. Since $0<v_1\le\log \scxo{D}$, the other two follow from those
for~$D$: $10^{6}v_1\le10^{6}\log \scxo{D}\le\log \scb{D}$ and $6\log v_1\le6\log\log \scxo{D}\le\log\log \scb{D}$. The quantities
listed do not involve the second scale.
\end{proof}

\begin{proof}[Proof of Theorem 2]
By Lemma~\ref{l:params}, \m{d\equiv7} modulo~$8$ and $N_0$ is prime to~$d$; also $1$ is prime to~$d$. So Input~8
applies with \m{M=N_0} and with \m{M=1}: the functions \m{\cE_1(v)=\R{d}{N_0}{\exp(v)}} and
\m{\cE_2(v)=\R{d}{1}{\exp(v)}} are polar, and Input~9 applies to them. Let $x$ be so large that
Input~7 gives admissible data~$D$ for~$P$ with first scale \m{\scx{D}=x} and the properties stated there, and that
\m{v_0=10^{6}\log x} is large enough for the conclusion of Input~9. Recall that $\scxo{D}$ and~$\scb{D}$ are the second and third
scale of~$D$ and $\Dm{D}$ its number of lifts. We show that \m{\R{d}{N_0}{x}\le x^{8\varepsilon+\frac12-1}}.

\emph{The window.} Put $X=(\None{D}/N_0)/\#\thick{D}$. The last property in Input~7, with
\m{\None{D}/N_0\ge1} and \m{\log \scb{D}>0}, shows that \m{\#\thick{D}>0} and that \m{0\le X\le\log \scb{D}}. By
Lemma~\ref{l:sizes}(a), $\log \scb{D}\le x^{1/100}\le\exp(v_0)$. So Input~9 gives a real number~$v_1$ with
\begin{equation}\label{e:window}
v_0\le v_1\le v_0^{2}\qquad\text{and}\qquad\R{d}{N_0}{\exp(v_1)}-X\cdot\R{d}{1}{\exp(v_1)}\;\le\;\exp(-v_0)=x^{-10^{6}}.
\end{equation}

\emph{The new data.} By Input~7, \m{v_0^{2}\le\log \scxo{D}}. So \m{10^{6}\log x=v_0\le v_1\le\log \scxo{D}}, and
Lemma~\ref{l:move} gives admissible data~$D'$ for~$P$ with second scale \m{\exp(v_1)}, which agree with~$D$ in
the first and third scale, the number of lifts, the thick places, the tops, and in $\fD{D}$, $\aD{D}$ and the
budget.

\emph{The identity.} Input~6 for~$D'$, with \m{\#\tops{D}=\None{D}/N_0} from Input~5, reads
\[
\rem{D'}=\aD{D}\cdot\Bigl(\#\thick{D}\cdot\bigl(\R{d}{N_0}{x}-\R{d}{N_0}{\exp(v_1)}\bigr)
+\frac{\None{D}}{N_0}\cdot\R{d}{1}{\exp(v_1)}\Bigr).
\]
Solving for \m{\R{d}{N_0}{x}}, which is possible because $\aD{D}$ and $\#\thick{D}$ are positive,
\begin{equation}\label{e:identity}
\R{d}{N_0}{x}=\frac{\Dm{D}\cdot\rem{D'}}{\Dm{D}\cdot\aD{D}\cdot\#\thick{D}}
+\Bigl(\R{d}{N_0}{\exp(v_1)}-X\cdot\R{d}{1}{\exp(v_1)}\Bigr).
\end{equation}

\emph{The two bounds.} The bracket in~\eqref{e:identity} is at most \m{x^{-10^{6}}} by~\eqref{e:window}.
Theorem~1 for~$D'$ gives \m{\Dm{D}\cdot\rem{D'}\le(\log \scb{D})^{2}\cdot\bud{D}=\fD{D}\cdot\aD{D}\cdot(\log \scb{D})^{5}}. So the
first term of~\eqref{e:identity} is at most
\[
\frac{\fD{D}\cdot(\log \scb{D})^{5}}{\Dm{D}\cdot\#\thick{D}}
\;\le\;\frac{\varphi(N_0)\,(\log \scb{D})^{6}}{\sqrt{x}}
\;=\;\varphi(N_0)\,x^{6\varepsilon-\frac12},
\]
where the inequality uses the value $\fD{D}=(\None{D}/N_0)\cdot\varphi(N_0)$ from Input~5, the bound
$\None{D}/N_0\le\log \scb{D}\cdot\#\thick{D}$ of Input~7, and \m{\Dm{D}\ge\sqrt{x}}; the equality is
\m{\log \scb{D}=x^{\varepsilon}}, again from Input~7.

\emph{The sizes.} By Input~5 and the conditions~(i),
$\varphi(N_0)\le\cP{P}\le10^{-12}\log\log \scb{D}=10^{-12}\,\varepsilon\log x$, and
\m{\varepsilon\log x\le\tfrac12\exp(2\varepsilon\log x)=\tfrac12x^{2\varepsilon}}. So the first term is at most
\m{\tfrac12x^{8\varepsilon-\frac12}}. As \m{x\ge2} by Lemma~\ref{l:sizes}(a) and \m{\varepsilon>0},
\m{x^{-10^{6}}\le\tfrac12x^{8\varepsilon-\frac12}}. Therefore
\m{\R{d}{N_0}{x}\le x^{8\varepsilon-\frac12}=x^{8\varepsilon+\frac12-1}}.
\end{proof}
\leanline{\lean{primeBound}}

The saving $x^{-1/2}$ comes from the division by~$\Dm{D}$: Theorem~1 bounds $\Dm{D}\cdot\rem{D'}$, not $\rem{D'}$, by a
quantity of the size of the trivial bound for the remainder, and $\Dm{D}$ is about $\sqrt{x}$. The number of lifts is
in this sense the source of the exponent~$\tfrac12$.

\sect{The zero-free half-plane}{s:grh}
\statement{inksoft}{\textbf{Input from another unit.}

\vspace{4pt}\phantomsection\label{i:10}\textbf{Input 10} (the one-sided criterion, after Landau; unit 151).\enspace
\emph{Let \m{M\ge1} be an integer, $c$ an element of \m{\ZZ/M} and $\beta$ a real number, with
\m{\tfrac12<\beta} and $c$ a unit of the ring \m{\ZZ/M}. Assume that for every prime number~$d$ whose class
modulo~$M$ is~$c$: \m{d\equiv7} modulo~$8$, $M$ is prime to~$d$, and \m{\R{d}{M}{x}\le x^{\beta-1}} for every
sufficiently large real number~$x$. Then for every complex number~$\rho$ with \m{\beta<\re\rho} and every Dirichlet
character~$\chi$ modulo~$M$, \m{\entireLrho{\chi}\ne0}.}}
\leanline{\lean{criterion}}

The criterion is read in the note on the endgame. Its proof takes one suitable prime~$d$ of the class; by Input~8
the Laplace transform of \m{v\mapsto\R{d}{M}{\exp(v)}} has a pole at every zero of an $L$-function modulo~$M$ that
is not cancelled by a zero of a twisted $L$-function, and Landau's theorem on integrals of functions of one sign
excludes poles to the right of~$\beta$ once the one-sided bound holds.

In the proof of Theorem~3 two facts of the library about characters are used. Let \m{M\ge1} be an integer, $\chi$
a Dirichlet character modulo~$M$, and $M'$ a multiple of~$M$ with \m{M'\ge1}; let $\chi'$ be the character
modulo~$M'$ induced by~$\chi$. First, $\chi'$ is principal if and only if $\chi$ is. Second, for every complex
number~$s$ such that $\chi$ is not principal or \m{s\ne1},
\begin{equation}\label{e:euler}
L(s,\chi')=L(s,\chi)\cdot\prod_{p\mid M'}\bigl(1-\chi(p)\,p^{-s}\bigr),
\end{equation}
the product over the primes~$p$ dividing~$M'$.

\begin{proof}[Proof of Theorem 3]
Let \m{M\ge1} be an integer, $\chi$ a Dirichlet character modulo~$M$ and $s$ a complex number with
\m{\re s>\tfrac12}, such that $\chi$ is not principal or \m{s\ne1}. Assume that \m{L(s,\chi)=0}; we derive a
contradiction. Put
\[
\varepsilon=\frac{\min(\re s,1)-\tfrac12}{51},\qquad\beta=8\varepsilon+\tfrac12 .
\]
Then \m{0<\varepsilon\le\tfrac1{102}<\tfrac1{100}} and \m{\tfrac12<\beta}, and
\m{\beta<\re s} because \m{8\varepsilon<\min(\re s,1)-\tfrac12}.

\emph{The modulus.} Let \m{N_0=18480\,M}, and let $\chi'$ be the character modulo~$N_0$ induced by~$\chi$. Then
\m{\entireLs{\chi'}=0}. Indeed, if $\chi$ is not principal, then $\chi'$ is not principal, and
\m{\entireLs{\chi'}=L(s,\chi')=0} by~\eqref{e:euler}. If $\chi$ is principal, then $\chi'$ is principal and
\m{s\ne1}, so \m{\entireLs{\chi'}=(s-1)L(s,\chi')=0}, again by~\eqref{e:euler}.

\emph{The class.} Let $c$ be the class of \m{9240\,M-1} modulo~$N_0$, that is, of \m{N_0/2-1}. It is a unit,
because \m{(9240\,M-1)^{2}=18480\,M\cdot(4620\,M-1)+1}. Let $d$ be a prime number in the class~$c$. Then
\m{P=(N_0,d)} are parameters: $18480$ divides~$N_0$; and \m{d+1\equiv N_0/2} modulo~$N_0$, so $N_0$ divides
\m{2(d+1)} and does not divide \m{d+1}. By Lemma~\ref{l:params}, \m{d\equiv7} modulo~$8$ and $N_0$ is prime to~$d$.
By Theorem~2 for $P$ and~$\varepsilon$,
\m{\R{d}{N_0}{x}\le x^{8\varepsilon+\frac12-1}=x^{\beta-1}} for every sufficiently large real number~$x$.

\emph{The criterion.} So the hypotheses of Input~10 hold for the modulus~$N_0$, the class~$c$ and~$\beta$. As
\m{\beta<\re s}, it gives \m{\entireLs{\chi'}\ne0}. This contradicts the first step.
\end{proof}
\leanline{\lean{grh}}

Theorem~2 is used here for every prime~$d$ of one residue class modulo~$N_0$, with one~$\varepsilon$ that depends
on~$s$; how large $x$ has to be may depend on~$d$. The modulus $18480\,M$ is chosen so that it can be the first
component of parameters, and the passage from $M$ to it costs nothing, because~\eqref{e:euler} is needed only in
one direction: a zero of \m{L(\,\cdot\,,\chi)} is a zero of \m{L(\,\cdot\,,\chi')}.

\sect{The critical line}{s:line}
\statement{inksoft}{\textbf{Input from another unit.}

\vspace{4pt}\phantomsection\label{i:11}\textbf{Input 11} (symmetry of the zeros; unit 106, background).\enspace
\emph{Let \m{M\ge1} be an integer, $\chi$ a Dirichlet character modulo~$M$, and $s$ a complex number with
\m{\re s\ne0}, \m{s\ne1} and \m{L(s,\chi)=0}. Then \m{\Gamma_{\chi}(s)=0} or \m{L(1-s,\chi^{-1})=0}.}}
\leanline{\lean{symm}}

Here \m{\Gamma_{\chi}} is the gamma factor of~$\chi$ in the library: \m{\Gamma_{\chi}(s)=\Gamma_{\RR}(s)} if
\m{\chi(-1)=1} and \m{\Gamma_{\chi}(s)=\Gamma_{\RR}(s+1)} otherwise, where
\m{\Gamma_{\RR}(s)=\pi^{-s/2}\,\Gamma(s/2)}. The library's function~$\Gamma$ takes the value~$0$ at its poles
\m{0,-1,-2,\dots}, so \m{\Gamma_{\chi}(s)=0} says that $s$ is a pole of the gamma factor. Input~11 is the
functional equation of the completed $L$-function, in the form that it takes for the zeros and for characters
that need not be primitive. It is a classical result; the catalogue of quoted results of this write-up records it
with its reference still to be confirmed. The library also provides the following facts, in which $\zeta^{*}$ is its
completed zeta function.
\begin{enumerate}[label=\textup{(Z\arabic*)},leftmargin=2.9em,itemsep=2pt]
\item For the Dirichlet character~$\chi$ modulo~$1$, \m{L(s,\chi)=\zeta(s)} for every complex number~$s$.
\item \m{\zeta(0)=-\tfrac12}.
\item For \m{s\ne0}, \m{\zeta(s)=\zeta^{*}(s)/\Gamma_{\RR}(s)}, where a quotient by~$0$ is~$0$.
\item \m{\zeta^{*}(1-s)=\zeta^{*}(s)} for every complex number~$s$.
\item \m{\Gamma_{\RR}(s)=0} if and only if \m{s=-2n} for a natural number~$n$; in particular
\m{\Gamma_{\RR}(s)\ne0} when \m{\re s>0}.
\end{enumerate}

\begin{proof}[Proof of Theorem 4]
Assume the half-plane statement. Let \m{M\ge1} be an integer, $\chi$ a Dirichlet character modulo~$M$ and $s$ a
complex number with \m{0<\re s<1} and \m{L(s,\chi)=0}. If \m{\re s>\tfrac12}, the half-plane statement gives
\m{L(s,\chi)\ne0}, since \m{s\ne1}; so \m{\re s\le\tfrac12}. Suppose that \m{\re s<\tfrac12}. Input~11 applies,
because \m{\re s\ne0} and \m{s\ne1}. The gamma factor does not vanish at~$s$, by (Z5): \m{\re s>0} and
\m{\re(s+1)>0}. Hence \m{L(1-s,\chi^{-1})=0}. But \m{\re(1-s)>\tfrac12} and \m{1-s\ne1}, so the half-plane
statement for the character~$\chi^{-1}$ gives \m{L(1-s,\chi^{-1})\ne0}, a contradiction. Therefore
\m{\re s=\tfrac12}.
\end{proof}
\leanline{\lean{zerosOnLine}}

\begin{proof}[Proof of Theorem 5]
Assume the half-plane statement, and let $s$ be a complex number with \m{\zeta(s)=0}, not of the form
\m{-2(n+1)} with $n$ a natural number, and with \m{s\ne1}. By (Z1), the half-plane statement for \m{M=1} says:
\m{\zeta(\rho)\ne0} for every complex number~$\rho$ with \m{\re\rho>\tfrac12} and \m{\rho\ne1}. With \m{\rho=s} this
gives \m{\re s\le\tfrac12}.

By (Z2), \m{s\ne0}. So $s$ is not of the form \m{-2n} with $n$ a natural number, and \m{\Gamma_{\RR}(s)\ne0} by
(Z5). By (Z3), \m{0=\zeta(s)=\zeta^{*}(s)/\Gamma_{\RR}(s)}, hence \m{\zeta^{*}(s)=0}, and \m{\zeta^{*}(1-s)=0} by
(Z4). As \m{1-s\ne0}, (Z3) gives \m{\zeta(1-s)=\zeta^{*}(1-s)/\Gamma_{\RR}(1-s)=0}. Now \m{1-s\ne1}; so if
\m{\re(1-s)>\tfrac12}, the half-plane statement with \m{\rho=1-s} would give \m{\zeta(1-s)\ne0}. Hence
\m{\re(1-s)\le\tfrac12}, that is, \m{\re s\ge\tfrac12}, and so \m{\re s=\tfrac12}.
\end{proof}
\leanline{\lean{zetaZerosOnLine}}

\sect{Where the results are used}{s:used}
Theorem~1 is used only in the proof of Theorem~2, and Theorem~2 only in the proof of Theorem~3. Theorems 3, 4
and~5 are the three inputs of the unit that holds the five final theorems: the third final theorem is Theorem~3,
the fourth is Theorem~4 applied to Theorem~3, and the first two are Theorem~5 applied to Theorem~3. Nothing else in
the artifact uses this unit.

\sect{Notes on the formal text}{s:notes}
\begin{enumerate}[label=\textup{\arabic*.},leftmargin=1.6em,itemsep=5pt]
\item \emph{Names.} The authors call Theorem~1 the inequality~$(\star)$ and Theorem~2 the one-sided bound. In the
formal text the parameters and the data of Theorem~1 are implicit arguments.
\item \emph{Theorem~2.} The two inequalities for~$\varepsilon$ form one hypothesis, a conjunction, written
\m{0<\varepsilon/6} and \m{\varepsilon<\tfrac1{100}}; the exponent is written \m{8\varepsilon+2^{-1}-1}. ``For
every sufficiently large real number'' renders Lean's ``eventually, as the number tends to infinity''. In
Definition~\ref{d:R}(d) the formal definition is made for every real~$x$; only large~$x$ matters. The constant
$\kap{d}{M}$ is the real part of the value at~$1$ of a function that the formal text defines for every complex
argument.
\item \emph{The record of the data.} The formal record stores the three logarithms, not the scales, and its
conditions~(i) are written with the library's logarithm, which is defined for all real numbers as the logarithm of
the absolute value. The conditions force \m{\log \scx{D}}, \m{\log \scxo{D}} and \m{\log \scb{D}} to be positive
(Lemma~\ref{l:sizes}(a) is proved from them in the order in which (i) lists them), so the iterated logarithms are
the usual ones. The record has one more component, a natural number that the formal text calls the exponent; a
condition makes it equal to \m{22\log_2\Dm{D}}. This write-up keeps the formal text's letter for it free for the
complex variable and does not name it: the degree of the tower, a function of the exponent in the formal text, is
\m{\None{D}/N_0} here, and the conditions on the exponent are those on \m{22\log_2\Dm{D}} in~(ii).
\item \emph{Input~7 is not used in full.} The proof of Theorem~2 does not use the equation
\m{\log \scxo{D}=x^{\varepsilon/6}}; it uses the inequality \m{(10^{6}\log x)^{2}\le\log \scxo{D}}, which leaves room for the
window.
\item \emph{The comparisons of sizes.} The formal text does not carry out the estimates of
Lemma~\ref{l:sizes}(c) and of the last step of the proof of Theorem~2. It proves them by a calculus of unit~152:
every quantity built from the data carries an upper or a lower bound that is the exponential of a real number,
found by Lean's search for instances, and an inequality between two quantities is reduced to an inequality
between two such exponents, which a decision procedure proves from the conditions~(i). The estimates in this note
are explicit, and they are mine; they use the same facts, which are the conditions (i) and~(ii), the definition
of~$\aD{D}$ and Input~5. Of Input~5 the formal proofs of this unit cite by name only the number of the tops; the
bounds for~$\cP{P}$ and the number of the infinite places enter through the calculus. All three statements were
read in the Lean text.
\item \emph{The new data.} Lemma~\ref{l:move} is not a declaration. The formal text builds $D'$ inside the proof of
Theorem~2 by replacing one component, and proves the conditions~(i) on the spot.
\item \emph{Lemma~\ref{l:params}.} The formal text cites two declarations of unit~152. The second says that the
order of the roots of unity, for every exponent of the tower, is prime to~$d$, and deduces this from the fact
that $\chi_K$ takes the value~$1$ at every prime dividing that order; the formal proofs of Theorems 2 and~3 use
it for the exponent~$0$ and keep the factor~$N_0$. The proof given here is direct.
\item \emph{The constant~$51$.} In the proof of Theorem~3 the formal text defines~$\varepsilon$ with the same
divisor~$51$; any divisor larger than~$50$ would do.
\item \emph{Theorems 4 and~5.} Each of the two formal proofs is one call of an automatic procedure that is given
the facts used above: Input~11 and (Z5) for Theorem~4, (Z1) to (Z5) for Theorem~5. The case distinctions written
out here are what that procedure has to find. They depend on the library's conventions, which are part of the
statements (Z3) and (Z5): a quotient by~$0$ is~$0$, and $\Gamma$ is~$0$ at its poles.
\item \emph{What the unit does not use.} Every declaration of the unit is used by the five final theorems.
\end{enumerate}

\subhead{How this note was written.} From the structured proofs of the unit: the program check of the dossier
covered \UnitSteps{} leaf steps of \UnitProofs{} proofs and compared the quotations of the imported declarations
with the source. In the Lean text I read the whole file of the unit, with the long comment at its head in which
the authors list the inputs, the scales and the errors. I read the statements of all eleven Inputs, and the
definitions behind every notion of §\ref{s:setting}: the parameters, the tower, the field, its places and their
classes, every field of the record of admissible data, the derived quantities, the weight, the remainder of the prime sum with
the functions it is built from, the prime sums and the remainder of the data, the polar functions, and the lines that
define the rows, the lifts, the index, the squared covolume and the main amount. I read the proof of Input~11, which
is short, and did not study the proofs of the other Inputs. In the library, at the pinned commit of its public repository, I read
the statements listed in the second table of Appendix~A. The letters and the forms of the notations are those of
the overview of this handout. The proofs follow the structured proofs, with these differences, which are mine:
Lemma~\ref{l:sizes} and the estimates at the end of the proof of Theorem~2 replace the calculus of sizes;
Lemma~\ref{l:params} is proved directly; the proofs of Theorems 4 and~5 are written out; and the explanations
between the proofs are mine.

\head{Appendix A. Concordance with the Lean artifact}
{\fontsize{8.6}{11.6}\selectfont\renewcommand{\leanfont}{\ttfamily\fontsize{7.7}{10}\selectfont\color{leancol}}
\setlength{\LTleft}{0pt}\setlength{\LTright}{\fill}\setlength{\tabcolsep}{0pt}
\begin{longtable}{@{}p{30mm}p{56mm}p{43mm}r@{}}
\textbf{In this note} & \textbf{Lean declaration} & \textbf{File} & \textbf{Lines}\\\midrule\endhead
\ConcordanceRows
\end{longtable}}

\subhead{The five statements of the unit as they are written in the artifact.} Theorems 1 to~5:
\quoted{star}
\quoted{primeBound}
\quoted{grh}
\quoted{zerosOnLine}
\quoted{zetaZerosOnLine}
Inputs 1 and~2:
\quoted{familyAB}
\quoted{longC}
Input~2 is stated in a part of its file where the rows, the two real numbers and its four hypotheses are declared
beforehand; the hypotheses are the arguments {\leanfont i}, {\leanfont hn}, {\leanfont I} and {\leanfont hV}.

\subhead{Legend.} For data {\leanfont D} of the formal text: {\leanfont D.logX}, {\leanfont D.logX₁},
{\leanfont D.logB} are \m{\log \scx{D}}, \m{\log \scxo{D}}, \m{\log \scb{D}}, and {\leanfont D.x}, {\leanfont D.x₁}, {\leanfont D.b}
are the scales; {\leanfont D.m} and {\leanfont D.r} are $\Dm{D}$ and~$\Dr{D}$; {\leanfont D.s} is the exponent
\m{22\log_2\Dm{D}}, and {\leanfont towerDeg D.s} is \m{\None{D}/N_0}; {\leanfont D.thick} is $\thick{D}$;
{\leanfont D.Φ} is the set of infinite places of~$\ED{D}$; {\leanfont P.N₀}, {\leanfont P.d} and {\leanfont P.c}
are $N_0$, $d$ and~$\cP{P}$. The class of a place is given by {\leanfont Parameters.placeMulEquiv} and, modulo a
divisor, by {\leanfont Parameters.placeClass}; the place \m{\nu\cdot\tau} is the product of
{\leanfont D.placeOfInt ν} and~$\tau$. The character \m{\xi\chi_K} is {\leanfont quadTwist ξ d}, and $\chi_K$ is
{\leanfont chiK d}; $\kap{d}{M}$ is the real part of {\leanfont logDerivMean d M 1}. In the formal proofs of this
unit the modulus of Theorem~3 is~{\leanfont q} and the complex number is~{\leanfont ρ}; in the file of the
statements they are {\leanfont N} and~{\leanfont s}. The window of Input~9 is written there with {\leanfont t},
{\leanfont u}, {\leanfont E} and {\leanfont E'} for $v_0$, $v_1$, $\cE_1$ and~$\cE_2$, and the class of Input~10 with
{\leanfont a} for~$c$. In the proof of Theorem~1 the formal text writes {\leanfont d} for \m{\rowDim{D}},
{\leanfont L} for \m{\log \scb{D}\cdot\bud{D}} and {\leanfont h} for~$\hA$.

\subhead{The facts of the library.} The library is Mathlib at the commit that the artifact pins in its manifest.

{\fontsize{8.6}{11.6}\selectfont\renewcommand{\leanfont}{\ttfamily\fontsize{7.7}{10}\selectfont\color{leancol}}
\setlength{\LTleft}{0pt}\setlength{\LTright}{\fill}\setlength{\tabcolsep}{0pt}
\begin{longtable}{@{}p{50mm}@{\hspace{4mm}}>{\raggedright\arraybackslash}p{86mm}@{}}
\textbf{In this note} & \textbf{Declaration of the library}\\\midrule\endhead
\mrow{induced character, principal or not}{DirichletCharacter.changeLevel, changeLevel\_eq\_one\_iff}
\mrow{the product formula \eqref{e:euler}}{DirichletCharacter.LFunction\_changeLevel}
\mrow{\m{\entireLs{\chi}} for principal~$\chi$}{DirichletCharacter.LFunctionTrivChar₁}
\mrow{the gamma factors \m{\Gamma_{\chi}} and \m{\Gamma_{\RR}}}{DirichletCharacter.gammaFactor, Complex.Gammaℝ}
\mrow{(Z1)}{DirichletCharacter.LFunction\_modOne\_eq}
\mrow{(Z2)}{riemannZeta\_zero}
\mrow{(Z3)}{riemannZeta\_def\_of\_ne\_zero}
\mrow{(Z4)}{completedRiemannZeta\_one\_sub}
\mrow{(Z5)}{Complex.Gammaℝ\_eq\_zero\_iff, Complex.Gammaℝ\_ne\_zero\_of\_re\_pos}
\end{longtable}}

\head{Notation}
\begingroup\fontsize{8.6}{10.9}\selectfont\raggedright
\setlength{\columnsep}{9mm}\setlength{\parskip}{0pt}
\begin{multicols}{2}
\nt{$P=(N_0,d)$}{parameters: $18480$ divides~$N_0$, $d$ is a prime \m{\equiv N_0/2-1} modulo~$N_0$}
\nt{$K$, $\chi_K$}{the field \m{\QQ(\sqrt{-d})}; the quadratic character of \m{\ZZ/d}, the Legendre symbol for odd~$d$}
\nt{$D$, $D'$}{admissible data for~$P$}
\nt{$\scx{D}$, $\scxo{D}$, $\scb{D}$}{the three scales of~$D$}
\nt{$x$}{a bare $x$ is the free real variable of Chapters I and II (Convention~0.3), as in the remainder $\R{d}{M}{x}$}
\nt{$\Dm{D}$, $\Dr{D}$}{the number of lifts and the rank of~$D$}
\nt{$\cP{P}$}{a constant attached to~$P$}
\nt{$\None{D}$, $\ED{D}$}{the order of the roots of unity; the field they generate over~$K$}
\nt{$\tau$, $\tau'$, $\cl{D}{\tau}$}{infinite places of~$\ED{D}$; the class in \m{(\ZZ/\None{D})^{\times}}}
\nt{$\thick{D}$, $\tops{D}$}{the thick places; the tops}
\nt{$\fD{D}$, $\aD{D}$}{the number of infinite places of~$\ED{D}$; the least multiple of $\None{D}$ that is
\m{\ge xb}}
\nt{$\bud{D}$}{the budget \m{\fD{D}\cdot\aD{D}\cdot(\log b)^3}}
\nt{$M$, $M'$}{moduli, integers \m{\ge1}}
\nt{$\nu$}{In Section~\ref{s:setting} and after it also an integer, never a character of the ideles (Convention~0.6)}
\nt{$n$, $p$}{a natural number; a prime number}
\nt{$\Lambda$, $\varphi$}{von Mangoldt's function; Euler's function}
\nt{$\lam{d}{M}{\nu}$}{the weight \m{\Lambda(\nu)(1-\chi_K(\nu))} on \m{\nu\equiv1} modulo~$M$}
\nt{$\chi$, $\chi'$, $\xi$}{Dirichlet characters; $\chi'$ induced by~$\chi$}
\nt{$s$, $\rho$}{complex numbers and not the exponent of the data (Convention~0.6)}
\nt{$L(s,\chi)$, $\entireLs{\chi}$}{the $L$-function; the entire function equal to it, or to
\m{(s-1)L(s,\chi)} for principal~$\chi$}
\nt{$\zeta$, $\zeta^{*}$}{the zeta function; the completed zeta function}
\nt{$\Gamma$, $\Gamma_{\RR}$, $\Gamma_{\chi}$, $\pi$}{the gamma function; \m{\pi^{-s/2}\Gamma(s/2)}; the gamma
factor of~$\chi$}
\nt{$\kap{d}{M}$, $\R{d}{M}{x}$}{a constant; the remainder of the weighted prime sum}
\nt{$\nomW{D}{\tau}$, $z$}{the nominal weight; a real number}
\nt{$\Ask{D}{z}{\tau}$, $\mainAsk{D}$}{a weighted prime sum; its expected value}
\nt{$\chowla{d}$}{Euler's constant minus \m{\re(L'/L)(1,\chi_K)}}
\nt{$\rem{D}$}{the remainder of the data~$D$}
\nt{$\rows{D}$, $\rowDim{D}$}{the space of rows; its dimension}
\nt{$\vartheta$, $\Lift{D}{\vartheta}$}{an $\Dm{D}$-fold twist; the composite of the lifts}
\nt{${F_1,\dots,F_n}$, $\sigma$}{algebraic modular forms; an automorphism of~$\CC$}
\nt{$\covol{F_1,\dots,F_n}{\sigma}$}{the squared covolume}
\nt{${y_1,\dots,y_n}$}{rows}
\nt{$\delta$, $V$}{real numbers: a bound for the index, the logarithm of a bound for the covolume}
\nt{$\ell$, $\eta_\ell$, $J$}{a natural number, mostly a prime; real numbers indexed by~$\ell$; a finite set of natural numbers (in the definition of the index)}
\nt{$\hmain{D}$, $h_{\mathrm{A}}$}{the main amount; the bound of Input~1 for one row, \m{\hmain{D}-m\cdot\rem{D}+\log b\cdot\bud{D}}}
\nt{$\varepsilon$, $\beta$}{real numbers; \m{\beta=8\varepsilon+\tfrac12} in §\ref{s:grh}}
\nt{$\cE$, $\cE_1$, $\cE_2$; $\omega$, $\omega_1$, $\omega_2$; $\Phi$}{functions \m{\RR\to\RR}; functions \m{\CC\to\CC}, residue functions; the transform, in the definition of polar functions}
\nt{$\res{d}{M}{\rho}$, $\ord_{\rho}$}{the residue of Input~8; the order of vanishing at~$\rho$}
\nt{$v$, $v_0$, $v_1$, $X$}{real numbers in the window of Input~9}
\nt{$c$}{a class modulo~$M$}
\end{multicols}
\endgroup

\end{document}
